\documentclass[journal]{IEEEtran}
\PassOptionsToPackage{x11names}{xcolor}

\usepackage{hyperref}
\usepackage{flushend}
\usepackage[inline]{enumitem}
\usepackage{array,multirow,graphicx}
\usepackage{float}
\usepackage{cite}
\usepackage{bm}
\usepackage{orcidlink}

\usepackage{xcolor}

\usepackage{footnote}
\makesavenoteenv{tabular}
\makesavenoteenv{table}
\makesavenoteenv{tabular*}
\makesavenoteenv{table*}

\usepackage{tikz}
\usepackage{pgfplots}
\usetikzlibrary{patterns}
    \pgfplotsset{compat=1.15}
\usepackage{colortbl}
\usepackage{graphicx}
\usepackage{multirow}
\usepackage{circuitikz}
\usepackage{dblfloatfix}
\usepackage[bottom]{footmisc}
\ifCLASSOPTIONcompsoc
    \usepackage[caption=false, font=normalsize, labelfont=sf, textfont=sf]{subfig}
\else
    \usepackage[caption=false, font=footnotesize]{subfig}
\fi

\usepackage{pifont}
\newcommand{\cmark}{\ding{51}}
\newcommand{\xmark}{\ding{55}}

\usepackage{amssymb}
\usepackage{amsmath}
\usepackage{mathtools}
\usepackage{booktabs}

\IEEEeqnarraydefcolsep{0}{\hoffset}
\newcommand*{\defeq}{\stackrel{\text{def}}{=}}

\newtheorem{defn}{Definition}[section]

\newcommand{\removable}[1]{{\color{red}#1}}
\newcommand{\suggested}[1]{{\color{teal}#1}}
\newcommand{\changed}[1]{{\color{orange}#1}}
\newcommand{\placeholder}[1]{\textcolor{blue}{[#1]}}

\hyphenation{op-tical net-works semi-conduc-tor}

\begin{document}

\title{Stochastic Harmonic Hosting Capacity \\ Considering Harmonic Recombination}
\author{Tom~Van~Acker~\orcidlink{0000-0001-6283-1457},%
        Jinwen~Xiang~\orcidlink{0000-0003-2107-0238},%
        and~Hakan~Ergun,~\IEEEmembership{Senior Member,~IEEE~\orcidlink{0000-0001-5171-1986}}%
\thanks{Tom Van Acker is with the Scenarios, Markets and Adequacy Department, Elia Transmission Belgium, 1000 Brussels, Belgium (e-mail: tom.vanacker@elia.com).}%
\thanks{Jinwen Xiang and Hakan Ergun are with the Research Group ELECTA, Department of Electrical Engineering, KU Leuven, 3000 Leuven, Belgium, and also with Etch-EnergyVille, 3600 Genk, Belgium (e-mail: jinwen.xiang@kuleuven.be, hakan.ergun@kuleuven.be).}
\thanks{This work is supported by the Energy Transition Fund by FOD Economy, Belgium, project HARMONIC. The General Board on Energy of the FOD Economy is not liable for any use of the presented work and subsequent results.}
\thanks{\textit{(Corresponding author: Tom~Van~Acker)}}}

\markboth{Journal of \LaTeX\ Class Files,~Vol.~14, No.~8, August~2015}%
{Shell \MakeLowercase{\textit{et al.}}: Bare Demo of IEEEtran.cls for IEEE Journals}

\maketitle

\begin{abstract}
The aggregation of harmonic currents from many power-electronic interfaced units depends on their phase angles, which are uncertain at the planning stage. When these angles differ, the harmonic currents of the units partly cancel, i.e., they recombine, instead of adding up arithmetically. This paper presents a stochastic harmonic hosting capacity method for phase-balanced three-phase power systems that captures harmonic recombination. The angle uncertainty is characterized through a harmonic state-space model of the converters and is represented in the hosting capacity optimization by polynomial chaos expansion, with one random variable per unit, and the angles are referenced to each unit's fundamental terminal voltage. Harmonic voltage and current limits are enforced through chance constraints. For fairness criteria that fix the ratios between the allocations of the units, the problem reduces to a small second-order cone program that equals the general formulation and recovers the deterministic hosting capacity exactly when the uncertainty vanishes. For the IEEE 30-bus system and the characteristic harmonics up to the 25th, recombination between units of a single converter technology, diode or thyristor bridges, grows with the harmonic order but remains limited, and the indicative summation exponents of IEC~61000-3-6 overestimate the hosting capacity by a factor of about two, also with underground cables. Cable resonances rotate the contributions of the units apart, which increases their cancellation. Aggregating the units at the mean of their harmonic current angles is unsafe, as it overestimates the hosting capacity wherever their contributions cancel, by more than a factor of two at the highest harmonics, whereas the presented chance-constrained approach guarantees that the harmonic limits are met with the prescribed probability.
\end{abstract}

\begin{IEEEkeywords}
Chance constraints, harmonic distortion, harmonic recombination, hosting capacity, polynomial chaos expansion.
\end{IEEEkeywords}

\IEEEpeerreviewmaketitle

\input{nomenclature}
\section{Introduction}
\label{sec:introduction}

\subsection{Background and Motivation}
\label{sec:background_and_motivation}

\IEEEPARstart{R}{eaching} a net-zero economy by 2050 requires extensive electrification of the energy system. During the past two decades, renewable generation, particularly wind and photovoltaic power, has grown steadily and is expected to accelerate towards 2050~\cite{international_energy_agency_world_2023}. At the same time, decarbonization targets require the electrification of transport, heating, and industrial processes, replacing fossil fuels with electricity. In Europe, studies project that the total electricity demand will triple by 2050~\cite{energyville_paths_2023}. Consequently, future power systems will be largely dominated by power-electronic interfaced resources on both the supply and demand sides. This transition fundamentally changes the harmonic behavior of power systems. Traditional grids, dominated by synchronous machines and mostly linear loads, exhibited limited harmonic distortion. Modern systems, however, increasingly rely on converters whose nonlinear characteristics generate a broad spectrum of harmonic currents. The aggregation of numerous distributed converters creates new challenges in maintaining harmonic voltage quality across transmission and distribution networks.

A key challenge arises from the stochastic nature of harmonic emissions. Measurements and laboratory studies show that harmonic magnitudes depend on loading, switching states, and internal control actions. More importantly, harmonic current phase angles, which govern the phasor summation of emissions from multiple sources, vary significantly. Even small angle variations can result in significant differences in harmonic voltage distortion, particularly in networks with many distributed converters. As the penetration of the converter increases, system operators need transparent metrics to assess whether additional units can be connected without exceeding harmonic limits. Harmonic hosting capacity provides such a metric by defining the maximum allowable harmonic emissions while ensuring compliance with power-quality standards. However, most existing methods rely on deterministic assumptions regarding harmonic magnitudes, phase angles, network parameters, and background distortion, leading to overly conservative or insufficiently protective results.

\subsection{Literature Review}
\label{sec:literature_review}

A growing body of literature presents approaches to determine the harmonic hosting capacity of a power system and to characterize uncertainty in harmonic power flow. Table~\ref{tab:literature_review} summarizes the publications discussed below, characterized by their mathematical method, network representation, stochastic method and tractability, and by which of the following features they represent stochastically:
\begin{enumerate*}
    \item [(a)] unit harmonic injection,
    \item [(b)] background harmonic distortion,
    \item [(c)] network parameters,
    \item [(d)] network exploitation, and
    \item [(e)] harmonic limits.
\end{enumerate*}

\begin{table*}[!t]
    \centering
    \caption{Literature review on harmonic hosting capacity and stochastic harmonic power flow, in order of discussion, enumerating the stochastic features~(a)--(e). The entry marked $\star$ is the deterministic predecessor of this work.}
    \label{tab:literature_review}
    \begin{tabular}{ c | c | c | c c | c | c c c c c | r r }
        \multirow{2}{*}{reference} & \multirow{2}{*}{objective} & \multirow{2}{*}{method$^{a}$} & \multicolumn{2}{c |}{network} & \multirow{2}{*}{stochastic$^{b}$} & \multicolumn{5}{c |}{stochastic features} & \multicolumn{2}{c}{tractability$^{c}$} \\
                                            &       &       & model & phases &       & (a)  & (b)   & (c)   & (d)   & (e)   & nodes & minutes \\ \hline
        \cite{noauthor_iec61000-3-6_2008}   & hhc   & ae    & nodal & bal.  & an.   &\cmark &\xmark &\xmark &\xmark &\cmark & --    & --    \\
        \cite{santos_considerations_2015}   & hhc   & ae    & nodal & bal.  & an.   &\cmark &\cmark &\xmark &\xmark &\xmark & --    & --    \\
        \cite{cho_analysis_2019}            & hhc   & pf    & full  & bal.  & mc    &\cmark &\xmark &\cmark &\xmark &\xmark & --    & --    \\ \arrayrulecolor{gray!20} \hline
        \cite{li2020}                       & hpf   & pf    & full  & bal.  & mc    &\cmark &\xmark &\cmark &\xmark &\xmark & 14    & --    \\
        \cite{xie_piecewise_2022}           & hpf   & pf    & full  & unb.  & an.   &\cmark &\xmark &\xmark &\xmark &\xmark & 8500  & 25    \\
        \cite{xie2022}                      & hpf   & pf    & full  & unb.  & an.   &\cmark &\xmark &\xmark &\xmark &\xmark & 8500  & 119   \\
        \cite{cadenazarate2024}             & hpf   & pf    & full  & bal.  & mc    &\cmark &\xmark &\xmark &\xmark &\xmark & --    & --    \\ \arrayrulecolor{gray!20} \hline
        \cite{sun_distribution_2012}        & hhc   & nl    & full  & bal.  & --    &\xmark &\xmark &\xmark &\xmark &\xmark & 5     & --    \\
        \cite{sun_incorporating_2012}       & hhc   & nl    & full  & bal.  & --    &\xmark &\xmark &\xmark &\xmark &\xmark & --    & --    \\
        \cite{sakar_integration_2018}       & hhc   & pf    & full  & bal.  & --    &\xmark &\xmark &\xmark &\xmark &\xmark & 3     & --    \\
        \cite{sun_distribution_2022}        & hhc   & cc    & full  & bal.  & scen. &\xmark &\cmark &\xmark &\cmark &\cmark & 16    & 30    \\
        \cite{nakhodchi2022}                & hhc   & ae    & nodal & bal.  & mc    &\cmark &\xmark &\xmark &\xmark &\cmark & --    & --    \\
        \cite{nakhodchi_impact_2023}        & hhc   & ae    & nodal & bal.  & mc    &\cmark &\xmark &\xmark &\xmark &\xmark & 64    & --    \\ \arrayrulecolor{black} \hline
        $\star$~\cite{vanacker_hosting_2026}& hhc   & soc   & full  & bal.  & --    &\xmark &\xmark &\xmark &\xmark &\xmark & 1880  & $<$\,6 \\ \arrayrulecolor{black} \hline
        This work                           & hhc   & qcqp/soc & full  & bal.  & pce   &\cmark &\xmark &\xmark &\xmark &\cmark & 30    & $<$\,0.04 \\
    \end{tabular}

    \vspace{2pt}
    \parbox{\textwidth}{\scriptsize
    $^{a}$Mathematical methods: (ae)~analytical expression, (pf)~power flow, (nl)~nonlinear optimal power flow, (soc)~second-order cone optimal power flow, (cc)~chance-constrained optimal power flow, (qcqp)~quadratically constrained quadratic program. Objective: (hhc)~harmonic hosting capacity, (hpf)~harmonic power flow. Phases: (bal.)~balanced, (unb.)~unbalanced.
    $^{b}$Stochastic methods: (mc)~Monte Carlo simulation, (an.)~analytical or moment-based, (scen.)~scenario-based, (pce)~polynomial chaos expansion.
    $^{c}$Tractability expressed as the number of nodes and the solve time in minutes.}
\end{table*}

\subsubsection{Standards and the summation law}

%The standard IEC~61000-3-6 allocates planning levels over installations connected to a point of common coupling and aggregates their emissions with a general summation law, $I_{\Harmonic} = (\sum_{\Unit} I_{\Unit,\Harmonic}^{\alpha})^{1/\alpha}$, with indicative exponents $\alpha = 1$ for $\Harmonic < 5$, $1.4$ for $5 \leq \Harmonic \leq 10$ and $2$ for $\Harmonic > 10$, and compares the 95\,\% distortion values with the planning levels~\cite{noauthor_iec61000-3-6_2008}. The summation law thus carries the diversity of the harmonic phase angles into the standard, but its exponents are indicative. Santos et al.\ used the summation law to determine the harmonic hosting capacity of transmission and distribution systems, represented the upstream network by a Norton equivalent with background distortion of unknown phase, and obtained a harmonic hosting capacity between 37\,MW and 160\,MW for the in-phase and anti-phase extremes; they also cautioned that the indicative exponents may under- or overestimate the aggregate of correlated sources such as wind and photovoltaic plants~\cite{santos_considerations_2015}. Field measurements confirm this caution: Wang et al.\ analyzed synchronized measurements in North-American residential distribution systems and found that low-order harmonic currents from different homes tend to add up, i.e., the effects is much less than claimed in the literature, while recombination occurs at high orders only~\cite{wang_characteristics_2017}. Cho et al.\ compared IEC~61000-3-6 with the deterministic current limits of IEEE Standard~519, set as a function of the short-circuit ratio~\cite{noauthor_ieee519_2022}, by Monte Carlo simulation on a multi-feeder distribution model, and showed that the latter can drive the voltage distortion of distribution systems up to twice the planning levels~\cite{cho_analysis_2019}. Trelles-Molina et al.\ extended the summation law from transmission systems to individual customer facilities with few harmonic sources and analyzed the probabilities that support the exponents~\cite{trellesmolina_extension_2022}. How much recombination a given network and converter population actually allow therefore remains an open question.
The standard IEC~61000-3-6 allocates planning levels over installations connected to a point of common coupling and aggregates their emissions, i.e., combines the harmonic currents of the individual installations into a single total current at the point of common coupling, with a general summation law, $I_{\Harmonic} = (\sum_{\Unit} I_{\Unit,\Harmonic}^{\alpha})^{1/\alpha}$, with indicative exponents $\alpha = 1$ for $\Harmonic < 5$, $1.4$ for $5 \leq \Harmonic \leq 10$ and $2$ for $\Harmonic > 10$, and compares the 95\,\% distortion values with the planning levels~\cite{noauthor_iec61000-3-6_2008}. The summation law thus carries the diversity of the harmonic phase angles into the standard, but its exponents are indicative. Santos et al.\ used the summation law to determine the harmonic hosting capacity of transmission and distribution systems, represented the upstream network by a Norton equivalent with background distortion of unknown phase, and obtained a harmonic hosting capacity between 37\,MW and 160\,MW for the in-phase and anti-phase extremes; they also cautioned that the indicative exponents may under- or overestimate the aggregate of correlated sources such as wind and photovoltaic plants~\cite{santos_considerations_2015}. Field measurements confirm this caution: Wang et al.\ analyzed synchronized measurements in North-American residential distribution systems and found that low-order harmonic currents from different homes tend to add up, i.e., the effects is much less than claimed in the literature, while recombination occurs at high orders only~\cite{wang_characteristics_2017}. Cho et al.\ compared IEC~61000-3-6 with the deterministic current limits of IEEE Standard~519, set as a function of the short-circuit ratio~\cite{noauthor_ieee519_2022}, by Monte Carlo simulation on a multi-feeder distribution model, and showed that the latter can drive the voltage distortion of distribution systems up to twice the planning levels~\cite{cho_analysis_2019}. Trelles-Molina et al.\ extended the summation law from transmission systems to individual customer facilities with few harmonic sources and analyzed the probabilities that support the exponents~\cite{trellesmolina_extension_2022}. How much recombination a given network and converter population actually allow therefore remains an open question.

%\subsubsection{Stochastic harmonic power flow}
\subsubsection{Probabilistic harmonic power flow}

Probabilistic harmonic power flow quantifies the uncertainty of the harmonic emissions and the associated network response. Li et al.\ separated uncertainty sources due to natural variability, i.e., aleatory, from those due to a lack of knowledge, i.e., epistemic, defined quantitative amplification indices as ratios of transfer to driving-point impedances, and assessed them with Monte Carlo simulation~\cite{li2020}. Monte Carlo and data-driven methods have since become the standard tool for probabilistic harmonic power flow: Xie and Sun proposed a piecewise probabilistic harmonic power flow for unbalanced residential distribution systems~\cite{xie_piecewise_2022}, Xie et al.\ a data-driven approach for distribution systems with photovoltaic generation~\cite{xie2022}, and Cadena-Zarate et al.\ a simulation-based approach for low-voltage networks with distributed energy resources~\cite{cadenazarate2024}. These methods quantify uncertainty, but do not optimize the hosting capacity.

\subsubsection{Stochastic harmonic hosting capacity}

Early harmonic hosting capacity methods incorporated harmonic limits within deterministic frameworks. Sun et al.\ incorporated harmonic voltage constraints into an optimal power flow for distribution network capacity assessment~\cite{sun_distribution_2012} and extended it to active networks~\cite{sun_incorporating_2012}, and Sakar et al.\ analyzed the hosting capacity of large photovoltaic plants in the presence of background distortion~\cite{sakar_integration_2018}. Sun and Harrison proposed the first optimization-based stochastic harmonic hosting capacity method, with percentile-based harmonic distortion limits as chance constraints within an ac optimal power flow, and showed that relaxing the risk level increases the admissible hosting capacity~\cite{sun_distribution_2022}. Nakhodchi et al.\ included the measured emission uncertainty of a fast charging station for electric vehicles through Bayesian statistics and the correlation between harmonic orders~\cite{nakhodchi2022}, and showed that the estimated hosting capacity depends strongly on how the medium-voltage network and the remote loads are modeled~\cite{nakhodchi_impact_2023}. The present authors determined an exact lower bound for the equitable harmonic hosting capacity with a second-order cone formulation of the full network, but with fixed harmonic current angles~\cite{vanacker_hosting_2026}, and studied the impact of underground cables on it~\cite{vanacker_impact_2025}.

\subsubsection{Research gaps}

Table~\ref{tab:literature_review} shows that no existing method simultaneously achieves:
\begin{enumerate*}
%    \item [(i)] tractable network-wide optimization,
    \item [(i)] computationally tractable network-wide optimization,
    \item [(ii)] explicit treatment of phase angle uncertainty, and
    \item [(iii)] chance-constraint enforcement of harmonic distortion limits.
\end{enumerate*}
Moreover, the recombination that a network and a converter population allow has not been quantified against the summation law of IEC~61000-3-6, and the conditions under which it occurs have not been identified.

\subsection{Contributions and Outline}
\label{sec:contribution_and_outline}

%In this paper, a holistic one-shot methodology is presented to determine the steady-state harmonic hosting capacity for an entire phase-balanced three-phase power system in a stochastic context. The presented method considers uncertainty on the angle of the injected harmonic current, enabling harmonic recombination, and limits the resulting individual and total harmonic voltage distortion, as well as the root-mean-square voltages and currents, using chance constraints. The problem is formulated from the perspective of the power system operator, who may anticipate the converter technologies its customers will connect, but not the operating point, parameters and control of the individual units. The angle of the harmonic current is therefore modeled as a random variable per unit, characterized for the anticipated technology, whereas its magnitude is the decision, and the harmonic hosting capacity is determined in terms of harmonic current magnitude using a fairness principle. The presented formulation uses the rectangular current-voltage variable space and incorporates:
In this paper, a holistic single-run methodology is presented to determine the steady-state harmonic hosting capacity for an entire phase-balanced three-phase power system in a stochastic context. The presented method considers uncertainty on the angle of the injected harmonic current, enabling harmonic recombination, and limits the resulting individual and total harmonic voltage distortion, as well as the root-mean-square voltages and currents, using chance constraints. The problem is formulated from the perspective of the power system operator, who may anticipate the converter technologies its customers will connect, but not the operating point, parameters and control of the individual units. The angle of the harmonic current is therefore modeled as a random variable per harmonic source unit, characterized for the anticipated technology, whereas its magnitude is the decision: the harmonic hosting capacity is determined as the harmonic current magnitude each unit may inject, using a fairness principle. The presented formulation uses the rectangular current-voltage variable space and incorporates:
\begin{itemize}
%    \item uncertainty on the angle of the injected harmonic current via polynomial chaos expansion,
    \item uncertainty on the angle of the injected harmonic current by applying the polynomial chaos expansion technique~\cite{sullivan_introduction_2015},
    \item exact representation of the harmonic impedance in the frequency domain and detailed treatment of transformers,
    \item enforcement of stochastic harmonic limits on bus voltage and component ampacity through chance constraints, and
    \item equitable distribution of the harmonic hosting capacity following a fairness principle.
\end{itemize}
%For fairness criteria that fix the ratios between the allocations of the units, e.g., absolute equality, the problem reduces to a second-order cone program with one pair of variables per harmonic. %
For fairness criteria that fix the ratios between the allocations of the units, e.g., absolute equality, which grants every unit the same harmonic current~\cite{vanacker_hosting_2026}, the problem reduces to a second-order cone program with one pair of variables per harmonic, i.e., a convex problem that standard solvers handle efficiently, as detailed in Section~\ref{sec:absolute_equality_socp}. %
The models and numerical illustrations are made freely available as part of the \textsc{Julia}-package \textsc{HarmonicPowerModels.jl}.\footnote{Available at \url{https://github.com/timmyfaraday/HarmonicPowerModels.jl}.}

The remainder of the paper is organized as follows: Section~\ref{sec:uncertainty} addresses the uncertainty of harmonic current phasors associated with power-electronic interfaced resources. Section~\ref{sec:mathematical_framework} presents the mathematical formulation of the stochastic harmonic hosting capacity problem. Section~\ref{sec:numerical_illustration} proves the equivalence of the presented methods on the IEEE 30-bus system and investigates the levers for harmonic recombination. Section~\ref{sec:conclusion} concludes the paper.

\section{Angle Uncertainty Related to Converters}
\label{sec:uncertainty}

\pgfplotsset{shhc/.style={
    axis lines=left, xmajorgrids=true, ymajorgrids=true, grid style={gray!10},
    label style={font=\footnotesize}, tick label style={font=\footnotesize},
    every axis plot/.append style={line width=0.6pt, mark size=1.5pt, black},
    legend style={draw=none, font=\scriptsize, legend columns=-1, at={(0.5,1.02)}, anchor=south},
    legend cell align=left},
    iec/.style={gray!60, line width=1.2pt, mark=none},
    sbase/.style={solid, mark=*},
    sbank/.style={densely dashed, mark=o},
    scable/.style={densely dotted, mark=square},
    scablec/.style={dashdotted, mark=triangle}}

This section addresses the uncertainty of harmonic current phase angles, related to the typical industrial load with diode-based or phase-controlled line-commutated converter.
% The steady-state harmonic state-space model of the line-commutated converter is established, which captures the mutual coupling among the harmonic phasors and their dependence on key passive parameters. The parameterized model is then integrated into a stochastic assessment framework. The key parameters are modeled as random variables with realistic bounds. Monte Carlo simulation is applied to generate probability distributions of harmonic phasors. Finally, the Gaussian copula correlation matrix of harmonic phase angles is obtained to indicate the inter-harmonic correlation, and the shared probability density functions are obtained.
The Monte Carlo simulation is applied to the steady-state harmonic state-space model of the converter to assess the uncertainty propagation from key circuit parameters to harmonic phase angles. The Gaussian copula correlation matrix of harmonic phase angles is obtained to indicate the inter-harmonic correlation, and the shared probability density functions are obtained.

\begin{figure}[!t]
\centering
\includegraphics[width=0.6\columnwidth]{figures/topo_indust_LCC.pdf}
\caption{Diode-based line-commutated converter with ac-side commutation inductance~$L_{ac}$, dc-side smoothing inductance~$L_{dc}$ and dc load.}
\label{fig:topo_indust_LCC} 
\end{figure}

\subsection{Steady-State Harmonic Model} 
% Internal frequency coupling is the key feature of converter-generated harmonics, which decide the recombination of harmonic phasors. Therefore, the accurate harmonic state-space model is built here to capture the nature commutation and harmonic coupling of line-commutated converter.
The steady-state harmonic state-space model is used here to accurately capture the nature commutation and internal harmonic coupling of the line-commutated converter. The diode-based converter shown in Fig. \ref{fig:topo_indust_LCC} is used to illustrate the modeling process. The converter periodically operates with sequential switching states, including 6 conducting states and 6 commutation states per fundamental period. Each state can be described by differential equations with time-varying coefficients that represent topology variations. The state equations for any switching state can be expressed as,
\begin{IEEEeqnarray}{l}
    (S_{dc} L_{ac} + L_{dc}) \frac{di_{dc}}{dt} = \nonumber \\
    S_{dc,ac}^{a} v_{ac}^a + S_{dc,ac}^{b} v_{ac}^b + S_{dc,ac}^{c} v_{ac}^c - R_{dc} i_{dc}, \label{eq:idc_t} \\
    % L_{ac} \frac{di_{ac}^a}{dt} - S_{dc,ac}^{a} L_{ac} \frac{di_{dc}}{dt} = \nonumber \\ S_{ac}^{a,a} v_{ac}^a + S_{ac}^{a,b} v_{ac}^b + S_{ac}^{a,c} v_{ac}^c, \label{eq:iac_t} \\
    L_{ac} \frac{di_{ac}^i}{dt} - S_{dc,ac}^{i} L_{ac} \frac{di_{dc}}{dt} = \nonumber \\ S_{ac}^{i,a} v_{ac}^a + S_{ac}^{i,b} v_{ac}^b + S_{ac}^{i,c} v_{ac}^c, \label{eq:iac_t}
\end{IEEEeqnarray}
where $v_{ac}^{i}$ and $i_{ac}^{i}$ are ac-side voltage and current, $v_{dc}$ and $i_{dc}$ are dc-side voltage and current, the superscript $i,j \in \{a, b, c\}$ represents the three phases, $S_{dc,ac}^{i}$, $S_{dc}$, and $S_{ac}^{i,j}$ are the piecewise switching functions reflecting the topology states.

The periodic equations in time-domain can be transformed into the harmonic state-space model in harmonic-domain through three operations:
(i) Fourier series expansion of state variables $x(t) \leftrightarrow \sum_{h=-N_h}^{N_h} X_h e^{jh\omega_1 t}$ with truncated order $N_h$,
(ii) frequency-domain differentiation $dx/dt \leftrightarrow \mathbf{Z}_{\Omega}\mathbf{X}$ where $\mathbf{Z}_{\Omega} = j\omega_1 \cdot \text{diag}(h)$, and
(iii) Toeplitz convolution matrices representing time-domain multiplication: $y(t) = S(t)x(t) \leftrightarrow \mathbf{Y} = \mathbf{S}_{conv}\mathbf{X}$.
The resulting unified system relates current harmonics $\mathbf{I} = [\mathbf{I}_{ac}^a, \mathbf{I}_{ac}^b, \mathbf{I}_{ac}^c]^T$ to voltage harmonics $\mathbf{U} = [\mathbf{V}_{ac}^a, \mathbf{V}_{ac}^b, \mathbf{V}_{ac}^c]^T$ through:

\begin{equation}
\mathbf{A}\mathbf{I} = \mathbf{B}\mathbf{U},
\label{eq:HSS}
\end{equation}
where matrices $\mathbf{A}$ and $\mathbf{B}$ contain Toeplitz convolution structures that explicitly capture frequency coupling. The model is parameterized by firing delay $\alpha_{\text{f}}$ (equal to $0$ for diode-based converter) and commutation overlap angle $\mu_{\text{c}}$, and inductances ($L_{ac}$, $L_{dc}$), providing a foundation for subsequent Monte Carlo simulation and uncertainty quantification.

\subsection{Methodology for Uncertainty Quantification}
In the planning stage, the uncertainty on key parameters comes from limited knowledge about the future integrated device, as well as manufacturing tolerances and operational variances. Therefore, inductance parameters $L_{ac}$ and $L_{dc}$ are modeled as random variables, and Monte Carlo simulation is used to assess the propagation of uncertainty.

\subsubsection{Realistic Parameter Distribution}
The uncertainties on $L_{ac}$ and $L_{dc}$ are modeled using truncated normal distributions:
\begin{equation}
L_{ac} \sim \mathcal{N}_{\text{trunc}}(\mu_{L_{ac}}, \sigma_{L_{ac}}^2, L_{ac,\min}, L_{ac,\max})
\label{truncated_normal_lac}
\end{equation}
\begin{equation}
L_{dc} \sim \mathcal{N}_{\text{trunc}}(\mu_{L_{dc}}, \sigma_{L_{dc}}^2, L_{dc,\min}, L_{dc,\max})
\label{truncated_normal_ldc}
\end{equation}
which are constrained within physically realistic bounds determined by engineering trade-offs at the industrial load design stage. For ac inductance $L_{ac}$, the upper bound $L_{ac,\max}$ is limited by the acceptable commutation angle and voltage drop, while the lower bound $L_{ac,\min}$ ensures adequate harmonic suppression and safe short-circuit current. For dc inductance $L_{dc}$, the upper bound $L_{dc,\max}$ is limited by dynamic response and size, while the lower bound $L_{dc,\min}$ ensures dc ripple to be within the specifications of power quality. The truncated normal distribution enforces physical bounds while maintaining mathematical tractability for Monte Carlo propagation.

% \begin{itemize}
%     \item \textit{ac-side inductance $L_{ac}$:} The upper bound $L_{ac,\max}$ is limited by acceptable commutation angle and voltage drop, while the lower bound $L_{ac,\min}$ ensures adequate harmonic suppression and safe short-circuit current.
%     \item \textit{dc-side inductance $L_{dc}$:} The upper bound $L_{dc,\max}$ is limited by required dynamic response, size, and cost, while the lower bound $L_{dc,\min}$ ensures dc ripple stays within specifications to prevent equipment aging and process quality degradation.
% \end{itemize}

\subsubsection{Monte Carlo Simulation}
Monte Carlo simulation is then implemented, by which 5000 pairs of random samples $(L_{ac}, L_{dc})$ are generated using rejection sampling to meet the defined constraints. For each of the parameter samples, a harmonic state-space model is built and solved. The iterative process involves: (i) recalculating the commutation angle $\mathsf{\mu}$ based on the sample pair of parameters; (ii) reconstructing the time-periodic switching functions $S_{dc,ac}^{i}(t)$, $S_{dc}(t)$, and $S_{ac}^{i,j}(t)$ using the updated commutation angle; and (iii) numerically solving the harmonic state-space equations to extract and store the magnitudes and phase angles of the ac-side harmonics.
% \begin{itemize}
%     \item Recalculating the commutation angle $\mathsf{\mu}$ based on the sampled pair of parameters.
%     \item Reconstructing time-periodic switching functions $S_{dc,ac}^{i}(t)$, $S_{dc}(t)$, and $S_{ac}^{i,j}(t)$ based on the updated commutation angle.
%     % \item Reconstructing time-periodic switching functions $\mathbf{S}_{dc,ac,conv}^{i}$, $\mathbf{S}_{dc,conv}$, and $\mathbf{S}_{ac,conv}^{i,j}$ based on the updated commutation angle.
%     \item Numerically solving the harmonic state-space equations to extract and store the magnitudes and phase angles of the ac-side harmonic currents of interest.
% \end{itemize}

\subsection{Dependency Analysis and Probability Density Function}
A case study is conducted based on a representative diode-based line-commutated converter unit, with nominal capacity of 1.25 MW and nominal voltage of 380 V, which is connected to the 35 kV bus through a rectifier transformer with a ratio of 35 kV/400 V. 
For the ac-side commutation inductance, the truncation bounds are $L_{ac,\min} = 0.01$ mH and $L_{ac,\max} = 0.03$ mH, corresponding to transformer short-circuit voltage ranging from $U_k = 5\%$ to $15\%$. For the dc-side smoothing inductance, the truncation bounds are $L_{dc,\min} = 0.05$ mH and $L_{dc,\max} = 0.35$ mH.

Fig.\ref{fig:harmonic_phasors_CI_LacLdc} shows the resulting phasor distributions for a subset of characteristic harmonics. In each harmonic phasor distribution plot, the black arrow represent the phasors with mean magnitude and mean phase angle, while the dark gray represent the phasors within the 95\% confidence intervals. The results indicate that (i) the harmonic phasors at different frequencies exhibit mutual dependency; (ii) harmonics at higher orders display progressively larger angle variations and uncertainties. 
% \begin{figure}[!t]
%     \centering
%     \includegraphics[width=\columnwidth]{figures/harmonic_phasors_CI_LacLdc.png}
%     \caption{Phasors of the characteristic harmonic currents of the converter of Fig.~\ref{fig:topo_indust_LCC} for 5000 Monte Carlo samples of~$L_{ac}$ and~$L_{dc}$: phasor with the mean magnitude and mean phase angle (red) and phasors within the 95\,\% confidence intervals (blue).}
%     \label{fig:harmonic_phasors_CI_LacLdc}
% \end{figure}

\usepgfplotslibrary{polar}
\providecommand{\Harmonic}{h}
\begin{figure}[!t]
    \centering
    % \includegraphics[width=0.4\textwidth]
    % {journals/shhc/data/phasor_ci/harmonic_phasors_diode.pdf}
    \pgfplotsset{phasorplot/.style={shhc, width=0.17\textwidth, ymin=0,
        xtick={0,90,180,270}, xticklabel={$\pgfmathprintnumber{\tick}^{\circ}$},
        yticklabel style={font=\tiny, anchor=north, inner sep=1.5pt, yshift=-2pt, /pgf/number format/1000 sep={}},
        every axis plot/.append style={line join=round}},
        outside/.style={only marks, mark=*, mark size=0.5pt, gray!35},
        inside/.style={only marks, mark=*, mark size=0.5pt, black!70},
        ci/.style={densely dashed, black, line width=0.5pt, mark=none},
        meanphasor/.style={-latex, black, line width=1pt, mark=none}}
    \subfloat[$\Harmonic = 1$]{%
    \begin{tikzpicture}
        \begin{polaraxis}[phasorplot, ymax=3200, ytick={2000}]
            \addplot[outside] table[x=angle, y=mag] {journals/shhc/data/phasor_ci/plot_outside_h1.dat};
            \addplot[inside] table[x=angle, y=mag] {journals/shhc/data/phasor_ci/plot_inside_h1.dat};
            \addplot[ci] table[x=angle, y=mag] {journals/shhc/data/phasor_ci/ci_h1.dat};
            \addplot[meanphasor] table[x=angle, y=mag] {journals/shhc/data/phasor_ci/mean_h1.dat};
        \end{polaraxis}
    \end{tikzpicture}}%
    \subfloat[$\Harmonic = 5$]{%
    \begin{tikzpicture}
        \begin{polaraxis}[phasorplot, ymax=700, ytick={300}]
            \addplot[outside] table[x=angle, y=mag] {journals/shhc/data/phasor_ci/plot_outside_h5.dat};
            \addplot[inside] table[x=angle, y=mag] {journals/shhc/data/phasor_ci/plot_inside_h5.dat};
            \addplot[ci] table[x=angle, y=mag] {journals/shhc/data/phasor_ci/ci_h5.dat};
            \addplot[meanphasor] table[x=angle, y=mag] {journals/shhc/data/phasor_ci/mean_h5.dat};
        \end{polaraxis}
    \end{tikzpicture}}%
    \subfloat[$\Harmonic = 7$]{%
    \begin{tikzpicture}
        \begin{polaraxis}[phasorplot, ymax=370, ytick={150}]
            \addplot[outside] table[x=angle, y=mag] {journals/shhc/data/phasor_ci/plot_outside_h7.dat};
            \addplot[inside] table[x=angle, y=mag] {journals/shhc/data/phasor_ci/plot_inside_h7.dat};
            \addplot[ci] table[x=angle, y=mag] {journals/shhc/data/phasor_ci/ci_h7.dat};
            \addplot[meanphasor] table[x=angle, y=mag] {journals/shhc/data/phasor_ci/mean_h7.dat};
        \end{polaraxis}
    \end{tikzpicture}}%
    \\
    \subfloat[$\Harmonic = 11$]{%
    \begin{tikzpicture}
        \begin{polaraxis}[phasorplot, ymax=230, ytick={100}]
            \addplot[outside] table[x=angle, y=mag] {journals/shhc/data/phasor_ci/plot_outside_h11.dat};
            \addplot[inside] table[x=angle, y=mag] {journals/shhc/data/phasor_ci/plot_inside_h11.dat};
            \addplot[ci] table[x=angle, y=mag] {journals/shhc/data/phasor_ci/ci_h11.dat};
            \addplot[meanphasor] table[x=angle, y=mag] {journals/shhc/data/phasor_ci/mean_h11.dat};
        \end{polaraxis}
    \end{tikzpicture}}%
    \subfloat[$\Harmonic = 13$]{%
    \begin{tikzpicture}
        \begin{polaraxis}[phasorplot, ymax=170, ytick={50}]
            \addplot[outside] table[x=angle, y=mag] {journals/shhc/data/phasor_ci/plot_outside_h13.dat};
            \addplot[inside] table[x=angle, y=mag] {journals/shhc/data/phasor_ci/plot_inside_h13.dat};
            \addplot[ci] table[x=angle, y=mag] {journals/shhc/data/phasor_ci/ci_h13.dat};
            \addplot[meanphasor] table[x=angle, y=mag] {journals/shhc/data/phasor_ci/mean_h13.dat};
        \end{polaraxis}
    \end{tikzpicture}}
    \caption{Phasors of a subset of characteristic harmonic currents of the diode-based converter (1000 of the 5000 Monte Carlo samples) with samples inside (dark gray) and outside (light gray) the 95\,\% confidence region of magnitude and phase angle (dashed), and the mean phasor (black arrow): (a) $\Harmonic = 1$; (b) $\Harmonic = 5$; (c) $\Harmonic = 7$; (d) $\Harmonic = 11$; (e) $\Harmonic = 13$.}
    \label{fig:harmonic_phasors_CI_LacLdc}
\end{figure}

% \usepgfplotslibrary{polar}
% \providecommand{\Harmonic}{h}
% \begin{figure*}[!t]
%     \centering
%     \includegraphics[width=0.7\textwidth]{journals/shhc/data/type2_phasor_ci/harmonic_phasors_pc.pdf}
%     % \pgfplotsset{phasorplotpc/.style={shhc, width=0.2\textwidth, ymin=0,
%     %     xtick={0,90,180,270}, xticklabel={$\pgfmathprintnumber{\tick}^{\circ}$},
%     %     yticklabel style={font=\tiny, anchor=north, inner sep=1.5pt, yshift=-2pt,
%     %         /pgf/number format/1000 sep={}},
%     %     every axis plot/.append style={line join=round}},
%     %     outsidepc/.style={only marks, mark=*, mark size=0.5pt, gray!35},
%     %     insidepc/.style={only marks, mark=*, mark size=0.5pt, black!70},
%     %     cipc/.style={densely dashed, black, line width=0.5pt, mark=none},
%     %     meanphasorpc/.style={-latex, black, line width=1pt, mark=none}}
%     % \subfloat[$\Harmonic = 1$]{%
%     % \begin{tikzpicture}
%     %     \begin{polaraxis}[phasorplotpc, ymax=3200, ytick={2000}]
%     %         \addplot[outsidepc] table[x=angle, y=mag] {journals/shhc/data/type2_phasor_ci/plot_outside_h1.dat};
%     %         \addplot[insidepc] table[x=angle, y=mag] {journals/shhc/data/type2_phasor_ci/plot_inside_h1.dat};
%     %         \addplot[cipc] table[x=angle, y=mag] {journals/shhc/data/type2_phasor_ci/ci_h1.dat};
%     %         \addplot[meanphasorpc] table[x=angle, y=mag] {journals/shhc/data/type2_phasor_ci/mean_h1.dat};
%     %     \end{polaraxis}
%     % \end{tikzpicture}}%
%     % \subfloat[$\Harmonic = 5$]{%
%     % \begin{tikzpicture}
%     %     \begin{polaraxis}[phasorplotpc, ymax=800, ytick={300}]
%     %         \addplot[outsidepc] table[x=angle, y=mag] {journals/shhc/data/type2_phasor_ci/plot_outside_h5.dat};
%     %         \addplot[insidepc] table[x=angle, y=mag] {journals/shhc/data/type2_phasor_ci/plot_inside_h5.dat};
%     %         \addplot[cipc] table[x=angle, y=mag] {journals/shhc/data/type2_phasor_ci/ci_h5.dat};
%     %         \addplot[meanphasorpc] table[x=angle, y=mag] {journals/shhc/data/type2_phasor_ci/mean_h5.dat};
%     %     \end{polaraxis}
%     % \end{tikzpicture}}%
%     % \subfloat[$\Harmonic = 7$]{%
%     % \begin{tikzpicture}
%     %     \begin{polaraxis}[phasorplotpc, ymax=370, ytick={150}]
%     %         \addplot[outsidepc] table[x=angle, y=mag] {journals/shhc/data/type2_phasor_ci/plot_outside_h7.dat};
%     %         \addplot[insidepc] table[x=angle, y=mag] {journals/shhc/data/type2_phasor_ci/plot_inside_h7.dat};
%     %         \addplot[cipc] table[x=angle, y=mag] {journals/shhc/data/type2_phasor_ci/ci_h7.dat};
%     %         \addplot[meanphasorpc] table[x=angle, y=mag] {journals/shhc/data/type2_phasor_ci/mean_h7.dat};
%     %     \end{polaraxis}
%     % \end{tikzpicture}}%
%     % \subfloat[$\Harmonic = 11$]{%
%     % \begin{tikzpicture}
%     %     \begin{polaraxis}[phasorplotpc, ymax=300, ytick={100}]
%     %         \addplot[outsidepc] table[x=angle, y=mag] {journals/shhc/data/type2_phasor_ci/plot_outside_h11.dat};
%     %         \addplot[insidepc] table[x=angle, y=mag] {journals/shhc/data/type2_phasor_ci/plot_inside_h11.dat};
%     %         \addplot[cipc] table[x=angle, y=mag] {journals/shhc/data/type2_phasor_ci/ci_h11.dat};
%     %         \addplot[meanphasorpc] table[x=angle, y=mag] {journals/shhc/data/type2_phasor_ci/mean_h11.dat};
%     %     \end{polaraxis}
%     % \end{tikzpicture}}%
%     % \subfloat[$\Harmonic = 13$]{%
%     % \begin{tikzpicture}
%     %     \begin{polaraxis}[phasorplotpc, ymax=200, ytick={50}]
%     %         \addplot[outsidepc] table[x=angle, y=mag] {journals/shhc/data/type2_phasor_ci/plot_outside_h13.dat};
%     %         \addplot[insidepc] table[x=angle, y=mag] {journals/shhc/data/type2_phasor_ci/plot_inside_h13.dat};
%     %         \addplot[cipc] table[x=angle, y=mag] {journals/shhc/data/type2_phasor_ci/ci_h13.dat};
%     %         \addplot[meanphasorpc] table[x=angle, y=mag] {journals/shhc/data/type2_phasor_ci/mean_h13.dat};
%     %     \end{polaraxis}
%     % \end{tikzpicture}}
%     \caption{Phasors of characteristic harmonic currents of the phase-controlled converter, showing 1000 of the 5000 Monte Carlo samples, with samples inside (dark gray) and outside (light gray) the 95\,\% confidence region of magnitude and phase angle (dashed), and the mean phasor (black arrow): (a) $\Harmonic = 1$; (b) $\Harmonic = 5$; (c) $\Harmonic = 7$; (d) $\Harmonic = 11$; (e) $\Harmonic = 13$. \placeholder{For visualization. To be commented out before submit.}}
%     \label{fig:harmonic_phasors_CI_LacLdc_pc}
% \end{figure*}

\subsubsection{Inter-Harmonic Dependency Analysis}
An important feature of the harmonics generated by the line-commutated converter is the inter-harmonic dependency due to frequency coupling, as indicated by the harmonic state-space model in ~\eqref{eq:HSS}. To characterize the dependencies among harmonic components, the correlation matrices with Kendall's Tau correlation coefficients are obtained.
% To assess this dependency, the correlation structure between different harmonic components is investigated. Correlation matrices, with Kendall's Tau correlation coefficients, are used to characterize the dependencies.
As shown by the phase angle correlation matrix $\mathbf{R}_{\text{phase}}$ in Fig.~\ref{fig:Correlation_Matrix_Copula}, phase angles have strong inter-harmonic dependency with the Gaussian copula parameters $\rho_{ij}^{Gauss} \geq 0.97$.
% indicating distinct characteristics in the inter-harmonic dependency structure for phase angles.

% \begin{figure}[!t]
%     \centering
%     \includegraphics[width=\columnwidth]{figures/Correlation_Matrix_Copula.png}
%     \caption{Gaussian copula correlation matrices of the harmonic current magnitudes and of the harmonic current phase angles across the harmonic orders, from 5000 Monte Carlo samples of~$L_{ac}$ and~$L_{dc}$.}
%     \label{fig:Correlation_Matrix_Copula}
% \end{figure}

% \providecommand{\Harmonic}{h}
% \begin{figure}[!t]
%     \centering
%     \pgfplotsset{corrplot/.style={shhc, grid=none, scale only axis, width=3.2cm, height=3.2cm,
%         xmin=-0.5, xmax=4.5, ymin=-0.5, ymax=4.5, y dir=reverse, axis on top,
%         xtick={0,...,4}, ytick={0,...,4}, xticklabels={1,5,7,11,13}, yticklabels={1,5,7,11,13},
%         tick style={draw=none}, xlabel={harmonic order $\Harmonic$},
%         colormap={gray}{color=(white) color=(gray!70)}, point meta min=0, point meta max=1,
%         nodes near coords={\pgfmathprintnumber[fixed, fixed zerofill, precision=2]{\pgfplotspointmeta}},
%         every node near coord/.append style={font=\tiny, black, anchor=center}},
%         corrcells/.style={matrix plot*, mesh/cols=5, point meta=explicit, draw=white, line width=0.4pt}}
%     % \subfloat[magnitude]{%
%     % \begin{tikzpicture}
%     %     \begin{axis}[corrplot, ylabel={harmonic order $\Harmonic$}]
%     %         \addplot[corrcells] table[x=col, y=row, meta=rho] {journals/shhc/data/copula_corr/corr_mag.dat};
%     %     \end{axis}
%     % \end{tikzpicture}}%
%     % \subfloat[phase angle]{%
%     \begin{tikzpicture}
%         \begin{axis}[corrplot]
%             \addplot[corrcells] table[x=col, y=row, meta=rho] {journals/shhc/data/copula_corr/corr_phase.dat};
%         \end{axis}
%     \end{tikzpicture}
%     % }
%     \caption{Dependence between harmonic current phase angles of the diode-based converter, Gaussian copula correlation matrices with parameters $\rho_{ij}^{Gauss}$ obtained from Kendall's rank correlation~$\tau$ of 5000 Monte Carlo samples, darker cells for stronger correlation.}
%     \label{fig:Correlation_Matrix_Copula}
% \end{figure}

\providecommand{\Harmonic}{h}
\begin{figure}[!t]
    \centering
    \pgfplotsset{corrplot/.style={shhc, grid=none, scale only axis, width=3.2cm, height=3.2cm,
        xmin=-0.5, xmax=4.5, ymin=-0.5, ymax=4.5, y dir=reverse, axis on top,
        xtick={0,...,4}, ytick={0,...,4}, xticklabels={1,5,7,11,13}, yticklabels={1,5,7,11,13},
        tick style={draw=none}, xlabel={harmonic order $\Harmonic$},
        colormap={gray}{color=(white) color=(gray!70)}, point meta min=0, point meta max=1,
        nodes near coords={\pgfmathprintnumber[fixed, fixed zerofill, precision=2]{\pgfplotspointmeta}},
        every node near coord/.append style={font=\tiny, black, anchor=center}},
        corrcells/.style={matrix plot*, mesh/cols=5, point meta=explicit, draw=white, line width=0.4pt}}
    % \subfloat[magnitude]{%
    % \begin{tikzpicture}
    %     \begin{axis}[corrplot, ylabel={harmonic order $\Harmonic$}]
    %         \addplot[corrcells] table[x=col, y=row, meta=rho] {journals/shhc/data/copula_corr/corr_mag.dat};
    %     \end{axis}
    % \end{tikzpicture}}%
    % \subfloat[phase angle]{%
    \begin{tikzpicture}
        % \begin{axis}[corrplot]
        %     \addplot[corrcells] table[x=col, y=row, meta=rho] {journals/shhc/data/copula_corr/corr_phase.dat};
        % \end{axis}
        \begin{axis}[corrplot, width=6.4cm, height=6.4cm, xmax=8.5, ymax=8.5,
            xtick={0,...,8}, ytick={0,...,8},
            xticklabels={1,5,7,11,13,17,19,23,25}, yticklabels={1,5,7,11,13,17,19,23,25},
            ylabel={harmonic order $\Harmonic$}]
            \addplot[corrcells, mesh/cols=9] table[x=col, y=row, meta=rho] {journals/shhc/data/copula_corr/corr_phase_to25_rho.dat};
        \end{axis}
    \end{tikzpicture}
    % }
    \caption{Dependence between harmonic current phase angles of the diode-based converter, Gaussian copula correlation matrices with parameters $\rho_{ij}^{Gauss}$ obtained from Kendall's rank correlation~$\tau$ of 5000 Monte Carlo samples, darker cells for stronger correlation.}
    \label{fig:Correlation_Matrix_Copula}
\end{figure}

% \subsubsection{Marginal Distribution Fitting}
% The marginal distributions of each harmonic component are modeled independently. The best-fit probability density function for each harmonic magnitude and phase angle is fitted to the data obtained from Monte Carlo simulation. The candidate distributions considered include Normal, Lognormal, Beta, Logistic and Exponential distributions.

% For each harmonic $h$, the marginal distribution $F_h(x_h)$ is estimated by fitting multiple parametric distributions and selecting the best-fit  distribution based on minimizing the root-mean-square error between the cumulative distribution functions simulated and fitted. The best-fit results are illustrated in Fig.\ref{fig:Bestfit_PDF}, showing the histograms and the fitted probability density functions of harmonic magnitudes and phase angles.

% Tables~\ref{tab:bestfit_phase} summarize the best-fit distributions and their estimated parameters for harmonic phase angles.
% \begin{table}[h] 
% \centering
% \caption{Best-Fit Distributions for Harmonic Phase Angles}
% \label{tab:bestfit_phase}
% \begin{tabular}{l l l l}
% \hline
% Order & distribution & error & parameters \\
% \hline
% 1st  & Beta    & $9.30 e^{-3}$ & $[ 10.8052, 20.2376, -107.9443, 13.0245]$ \\
% 5th  & Beta    & $1.97 e^{-3}$ & $[11.8144, 23.6725, -2.6825, 70.8689]$ \\
% 7th  & Beta    & $1.54 e^{-3}$ & $[ 9.7918, 18.8775, -29.9737, 88.2789]$ \\
% 11th & Beta    & $7.41 e^{-4}$ & $[ 12.4463, 22.1493, 59.5585, 164.3101]$ \\
% 13th & Beta    & $6.84 e^{-4}$ & $[ 10.4088, 17.2887, 29.4667, 172.4882]$ \\
% \hline
% \end{tabular}
% \end{table}

% The parameters of the Beta distribution are given as $[\alpha, \beta, \ell, \theta]$, where $\alpha$ and $\beta$ are the shape parameters, and $\ell$ and $\theta$ are the location and scale parameters, respectively.

% \begin{figure*}[h!]
%     \centering
%     \includegraphics[width=0.8\textwidth]{figures/Bestfit_PDF.png}
%     \caption{Histograms of the harmonic current magnitudes and phase angles from 5000 Monte Carlo samples of~$L_{ac}$ and~$L_{dc}$, with the best-fit probability density functions of Table~\ref{tab:bestfit_phase}.}
%     \label{fig:Bestfit_PDF} 
% \end{figure*}

\subsubsection{Shared Beta Distribution for Phase Angles}
Since strong inter-harmonic dependencies are proved for phase angles, a shared Beta distribution is used to represent the normalized phase angle distributions across the harmonic orders. In this way, only two shape parameters, $\alpha$ and $\beta$, are estimated for all harmonics, while each harmonic retains its own location and scale parameters to capture the specific range of phase angles.

Firstly, a normalization procedure is applied to the unwrapped phase angle data of each harmonic to map them to the unit interval $[0,1]$. Let $\phi_h = \{\phi_{h,i}\}_{i=1}^{N}$ denote the unwrapped phase angle samples in degrees for harmonic $h \in \{1, 5, 7, 11, 13, 17, 19, 23, 25\}$, where unwrapping avoids discontinuities. For each harmonic, the observed min–max values are $v_{\min,h} = \min(\phi_h)$ and $v_{\max,h} = \max(\phi_h)$.
% \begin{equation}
% v_{\min,h} = \min(\phi_h), \qquad v_{\max,h} = \max(\phi_h).
% \end{equation}
The sampled data of each harmonic phase angle are then normalized to the unit interval via the min-max transformation:
\begin{equation}
u_{h,i} = \frac{\phi_{h,i} - v_{\min,h}}{v_{\max,h} - v_{\min,h}} \in (0,1).
\end{equation}

Secondly, all normalized samples from the five harmonics are pooled into a single dataset:
\begin{equation}
\mathcal{U} = \bigcup_{h} \{u_{h,i}\}_{i=1}^{N}.
\end{equation}

% A single Beta distribution with estimated shape parameters $\hat{\alpha}$ and $\hat{\beta}$ is fitted to $\mathcal{U}$ via maximum likelihood estimation, with location and scale fixed to the standard interval $[0,1]$:
% \begin{equation}
% u \sim \mathrm{Beta}(\hat{\alpha}, \hat{\beta}).
% \end{equation}

% The results of the fitted parameters are $\hat{\alpha} = 4.0414$, $\hat{\beta} = 4.7762$. The Beta probability density function is given by:
% \begin{equation}
% f_{\mathrm{Beta}}(u; \hat{\alpha}, \hat{\beta}) = \frac{u^{\hat{\alpha}-1}(1-u)^{\hat{\beta}-1}}{\int_0^1 t^{\hat{\alpha}-1}(1-t)^{\hat{\beta}-1}\,dt}, \quad u \in (0,1).
% \end{equation}

A single Beta distribution $f_{\mathrm{Beta}}(u; \hat{\alpha}, \hat{\beta})$ with estimated shape parameters $\hat{\alpha}$ and $\hat{\beta}$ is fitted to $\mathcal{U}$ via maximum likelihood estimation, with location and scale fixed to the standard interval $[0,1]$.

Thirdly, the fitted shared Beta distribution is transformed back to each harmonic's phase angle scale to assess the goodness-of-fit.
% For each harmonic $h$, the denormalization is performed as:
% \begin{equation}
% \label{eq:beta_denorm}
% f_{\phi_h}(\phi) = \frac{1}{v_{\max,h} - v_{\min,h}} \, f_{\mathrm{Beta}}\!\left(\frac{\phi - v_{\min,h}}{v_{\max,h} - v_{\min,h}}; \hat{\alpha}, \hat{\beta}\right),
% \end{equation}
% where $\phi\in [v_{\min,h}, v_{\max,h}]$ is the original phase angle range for harmonic $h$.

The comparison plots between the empirical histograms and the denormalized shared Beta probability density functions for each harmonic phase angle are shown in Fig.~\ref{fig:Shared_Beta_Fit_diode}.
Table~\ref{tab:beta_phase_params} lists the shared shape parameters, harmonic-specific locations $v_{\min,h}$ and scales $(v_{\max,h} - v_{\min,h})$, and the root-mean-square errors between the empirical histograms from the Monte Carlo simulation and the denormalized Beta probability density functions. Both the visual comparison and the low root-mean-square errors indicate a good fit of the shared Beta distribution to the simulated data across all harmonics.
% The goodness-of-fit of the shared Beta distribution model is evaluated for each harmonic phase angle by comparing the denormalized Beta probability density function against the empirical histograms from Monte Carlo simulation. The root-mean-square error is computed between the empirical histogram and the denormalized Beta probability density function, with the results summarized in last column of Table~\ref{tab:beta_phase_params}.
% \begin{figure}[!b]
%     \centering
%     % \includegraphics[width=0.8\textwidth]{figures/Shared_Beta_Fit.png}
%     \includegraphics[width=\columnwidth]{figures/Shared_Beta_Fit.png}
%     \caption{Histograms of the harmonic current phase angles from 5000 Monte Carlo samples of~$L_{ac}$ and~$L_{dc}$, and the denormalized shared Beta probability density functions with parameters listed in Table~\ref{tab:beta_phase_params}.}
%     \label{fig:Shared_Beta_Fit} 
% \end{figure}

% \begin{figure*}[t]
%     \centering
%     \pgfplotsset{phaseplot/.style={shhc, width=0.2\textwidth, height=3cm, ymin=0,
%         xlabel={phase [deg]}, scaled y ticks=false, yticklabel style={/pgf/number format/fixed},
%         every axis plot/.append style={mark=none}},
%         mcs/.style={ybar interval, fill=gray!30, draw=gray!70, line width=0.2pt},
%         beta/.style={solid, black, line width=0.9pt}}
%     \subfloat[$h = 1$]{%
%     \begin{tikzpicture}
%         \begin{axis}[phaseplot, xtick={-106,-103,-100}, ylabel style={align=center}, ylabel={probability\\density\\{[1/deg]}}]
%             \addplot[mcs] table[x=x, y=density] {journals/shhc/data/shared_beta_phase/hist_interval_h1.dat};
%             \addplot[beta] table[x=x, y=pdf] {journals/shhc/data/shared_beta_phase/pdf_h1.dat};
%         \end{axis}
%     \end{tikzpicture}}%
%     \subfloat[$h = 5$]{%
%     \begin{tikzpicture}
%         \begin{axis}[phaseplot, xtick={10,20,30}]
%             \addplot[mcs] table[x=x, y=density] {journals/shhc/data/shared_beta_phase/hist_interval_h5.dat};
%             \addplot[beta] table[x=x, y=pdf] {journals/shhc/data/shared_beta_phase/pdf_h5.dat};
%         \end{axis}
%     \end{tikzpicture}}%
%     \subfloat[$h = 7$]{%
%     \begin{tikzpicture}
%         \begin{axis}[phaseplot, xtick={-20,0,20}]
%             \addplot[mcs] table[x=x, y=density] {journals/shhc/data/shared_beta_phase/hist_interval_h7.dat};
%             \addplot[beta] table[x=x, y=pdf] {journals/shhc/data/shared_beta_phase/pdf_h7.dat};
%         \end{axis}
%     \end{tikzpicture}}%
%     \subfloat[$h = 11$]{%
%     \begin{tikzpicture}
%         \begin{axis}[phaseplot, xtick={90,120,150}]
%             \addplot[mcs] table[x=x, y=density] {journals/shhc/data/shared_beta_phase/hist_interval_h11.dat};
%             \addplot[beta] table[x=x, y=pdf] {journals/shhc/data/shared_beta_phase/pdf_h11.dat};
%         \end{axis}
%     \end{tikzpicture}}%
%     \subfloat[$h = 13$]{%
%     \begin{tikzpicture}
%         \begin{axis}[phaseplot, xtick={60,100,140}]
%             \addplot[mcs] table[x=x, y=density] {journals/shhc/data/shared_beta_phase/hist_interval_h13.dat};
%             \addplot[beta] table[x=x, y=pdf] {journals/shhc/data/shared_beta_phase/pdf_h13.dat};
%         \end{axis}
%     \end{tikzpicture}}
%     \caption{Distribution of harmonic current phase angles of the diode-based converter, histograms of 5000 Monte Carlo samples (gray, 50 bins) against the denormalized shared Beta probability density functions (black) with parameters listed in Table~\ref{tab:beta_phase_params}: (a) $h = 1$; (b) $h = 5$; (c) $h = 7$; (d) $h = 11$; (e) $h = 13$.}
%     \label{fig:Shared_Beta_Fit}
% \end{figure*}

\providecommand{\diodedatatoXXV}{journals/shhc/data/shared_beta_phase/}
\begin{figure}[!b]
    \centering
    \pgfplotsset{diodephaseXXV/.style={shhc, scale only axis, width=0.20\columnwidth, height=1.3cm, ymin=0,
        scaled y ticks=false, scaled x ticks=false,
        yticklabel style={/pgf/number format/fixed, /pgf/number format/precision=3},
        tick label style={font=\tiny}, label style={font=\scriptsize},
        max space between ticks=20pt,
        every axis plot/.append style={mark=none}},
        diodeleftXXV/.style={ylabel={density [1/deg]}},
        diodebottomXXV/.style={xlabel={phase [deg]}},
        diodemcsXXV/.style={ybar interval, fill=gray!30, draw=gray!70, line width=0.2pt},
        diodebetaXXV/.style={solid, black, line width=0.7pt}}
    \subfloat[$h=1$]{%
    \begin{tikzpicture}
        \begin{axis}[diodephaseXXV, xtick={-106,-103,-100}, diodeleftXXV]
            \addplot[diodemcsXXV] table[x=x,y=density] {\diodedatatoXXV/hist_interval_h1.dat};
            \addplot[diodebetaXXV] table[x=x,y=pdf] {\diodedatatoXXV/pdf_h1.dat};
        \end{axis}
    \end{tikzpicture}}%
    \subfloat[$h=5$]{%
    \begin{tikzpicture}
        \begin{axis}[diodephaseXXV]
            \addplot[diodemcsXXV] table[x=x,y=density] {\diodedatatoXXV/hist_interval_h5.dat};
            \addplot[diodebetaXXV] table[x=x,y=pdf] {\diodedatatoXXV/pdf_h5.dat};
        \end{axis}
    \end{tikzpicture}}%
    \subfloat[$h=7$]{%
    \begin{tikzpicture}
        \begin{axis}[diodephaseXXV]
            \addplot[diodemcsXXV] table[x=x,y=density] {\diodedatatoXXV/hist_interval_h7.dat};
            \addplot[diodebetaXXV] table[x=x,y=pdf] {\diodedatatoXXV/pdf_h7.dat};
        \end{axis}
    \end{tikzpicture}}\par
    \subfloat[$h=11$]{%
    \begin{tikzpicture}
        \begin{axis}[diodephaseXXV, xtick={80,120,160}, diodeleftXXV]
            \addplot[diodemcsXXV] table[x=x,y=density] {\diodedatatoXXV/hist_interval_h11.dat};
            \addplot[diodebetaXXV] table[x=x,y=pdf] {\diodedatatoXXV/pdf_h11.dat};
        \end{axis}
    \end{tikzpicture}}%
    \subfloat[$h=13$]{%
    \begin{tikzpicture}
        \begin{axis}[diodephaseXXV]
            \addplot[diodemcsXXV] table[x=x,y=density] {\diodedatatoXXV/hist_interval_h13.dat};
            \addplot[diodebetaXXV] table[x=x,y=pdf] {\diodedatatoXXV/pdf_h13.dat};
        \end{axis}
    \end{tikzpicture}}%
    \subfloat[$h=17$]{%
    \begin{tikzpicture}
        \begin{axis}[diodephaseXXV, xtick={-200,-100}]
            \addplot[diodemcsXXV] table[x=x,y=density] {\diodedatatoXXV/hist_interval_h17.dat};
            \addplot[diodebetaXXV] table[x=x,y=pdf] {\diodedatatoXXV/pdf_h17.dat};
        \end{axis}
    \end{tikzpicture}}\par
    \subfloat[$h=19$]{%
    \begin{tikzpicture}
        \begin{axis}[diodephaseXXV, diodeleftXXV, diodebottomXXV]
            \addplot[diodemcsXXV] table[x=x,y=density] {\diodedatatoXXV/hist_interval_h19.dat};
            \addplot[diodebetaXXV] table[x=x,y=pdf] {\diodedatatoXXV/pdf_h19.dat};
        \end{axis}
    \end{tikzpicture}}%
    \subfloat[$h=23$]{%
    \begin{tikzpicture}
        \begin{axis}[diodephaseXXV, diodebottomXXV]
            \addplot[diodemcsXXV] table[x=x,y=density] {\diodedatatoXXV/hist_interval_h23.dat};
            \addplot[diodebetaXXV] table[x=x,y=pdf] {\diodedatatoXXV/pdf_h23.dat};
        \end{axis}
    \end{tikzpicture}}%
    \subfloat[$h=25$]{%
    \begin{tikzpicture}
        \begin{axis}[diodephaseXXV, diodebottomXXV]
            \addplot[diodemcsXXV] table[x=x,y=density] {\diodedatatoXXV/hist_interval_h25.dat};
            \addplot[diodebetaXXV] table[x=x,y=pdf] {\diodedatatoXXV/pdf_h25.dat};
        \end{axis}
    \end{tikzpicture}}
    \caption{Harmonic current phase angles of the diode-based converter, histograms of 5000 Monte Carlo samples per harmonic (gray, 50 bins) against denormalized shared Beta densities (black), with parameters listed in Table~\ref{tab:beta_phase_params}: (a) $h=1$; (b) $h=5$; (c) $h=7$; (d) $h=11$; (e) $h=13$; (f) $h=17$; (g) $h=19$; (h) $h=23$; (i) $h=25$.}
    \label{fig:Shared_Beta_Fit_diode}
\end{figure}

% \begin{table}[!t]
%     \centering
%     \caption{Shared Beta distribution parameters and fit quality for harmonic phase angles of the diode-based converter}
%     \label{tab:beta_phase_params}
%     \scriptsize
%     \begin{tabular}{l c c c c}
%     Order & $v_{\min}$ [deg] & $v_{\max}$ [deg] & $v_{\max} - v_{\min}$ [deg] & error \\
%     \hline
%     1st  & $-106.70$ & $-99.43$   & 7.27   & 0.0226 \\
%     5th  & 4.75      & 40.05      & 35.30  & 0.0037 \\
%     7th  & $-22.77$  & 27.51      & 50.28  & 0.0029 \\
%     11th & 79.69     & 164.01     & 84.32  & 0.0014 \\
%     13th & 49.38     & 147.67     & 98.29  & 0.0011 \\
%     \end{tabular}
    
%     \vspace{2pt}
%     \parbox{\columnwidth}{\scriptsize All harmonics share the same shape parameters ($\hat{\alpha} = 4.041$, $\hat{\beta} = 4.776$).\\
%     Location parameters $v_{\min}$ and scale parameter $(v_{\max} - v_{\min})$ are harmonic-specific.}
% \end{table}

\begin{table}[!b]
    \centering
    \caption{Shared Beta distribution parameters and fit quality for harmonic current angles of diode-based converter}
    \label{tab:beta_phase_params}
    \scriptsize
    \setlength{\tabcolsep}{3pt}
    \begin{tabular}{l c c c c c}
    Order & $v_{\min}$ [deg] & $v_{\max}$ [deg] & $v_{\max}-v_{\min}$ [deg] & Mean [deg] & error \\
    \hline
    1st & -106.43 & -99.26 & 7.17 & -103.42 & 0.0403 \\
    5th & 4.80 & 40.90 & 36.10 & 20.87 & 0.0055 \\
    7th & -22.52 & 27.18 & 49.71 & 0.08 & 0.0039 \\
    11th & 77.87 & 163.96 & 86.09 & 118.52 & 0.0023 \\
    13th & 48.98 & 148.35 & 99.37 & 94.10 & 0.0016 \\
    17th & -239.51 & -76.70 & 162.81 & -157.65 & 0.0020 \\
    19th & 80.08 & 266.23 & 186.15 & 170.26 & 0.0013 \\
    23th & -209.79 & 34.31 & 244.10 & -100.85 & 0.0005 \\
    25th & -246.50 & 15.56 & 262.06 & -138.77 & 0.0011 \\
    \end{tabular}
    \vspace{2pt}
    \parbox{\columnwidth}{\scriptsize All harmonics share the jointly fitted shape parameters $\hat{\alpha}=3.731$ and $\hat{\beta}=4.439$. \\
    Location parameters $v_{\min}$ and scale parameter $(v_{\max} - v_{\min})$ are harmonic-specific.}
\end{table}

\begin{table}[!b]
    \centering
    \caption{Shared Beta distribution parameters and fit quality for harmonic current angles of phase-controlled converter}
    \label{tab:beta_phase_params_pc}
    \scriptsize
    \setlength{\tabcolsep}{3pt}
    \begin{tabular}{l c c c c c}
    Order & $v_{\min}$ [deg] & $v_{\max}$ [deg] & $v_{\max}-v_{\min}$ [deg] & Mean [deg] & error \\
    \hline
    1st & -136.66 & -110.06 & 26.60 & -122.56 & 0.0069 \\
    5th & -143.07 & -9.30 & 133.77 & -71.00 & 0.0014 \\
    7th & -263.84 & -52.49 & 211.36 & -143.42 & 0.0012 \\
    11th & -248.03 & 46.71 & 294.74 & -88.68 & 0.0005 \\
    13th & -353.10 & 5.22 & 358.32 & -156.41 & 0.0004 \\
    17th & -346.15 & 104.27 & 450.42 & -103.60 & 0.0003 \\
    19th & -447.14 & 62.43 & 509.57 & -169.55 & 0.0003 \\
    23th & -442.00 & 161.32 & 603.31 & -116.36 & 0.0002 \\
    25th & -180.49 & 479.68 & 660.17 & 178.94 & 0.0002 \\
    \end{tabular}
    \vspace{2pt}
    \parbox{\columnwidth}{\scriptsize All harmonics share the jointly fitted shape parameters $\hat{\alpha}=4.1612$ and $\hat{\beta}=3.5541$. \\
    Location parameters $v_{\min}$ and scale parameter $(v_{\max} - v_{\min})$ are harmonic-specific.
    % Angles are unwrapped about their respective circular means.
    }
\end{table}

\subsection{Second Converter Type: Phase-Controlled Rectifier}
\label{sec:converter_types}
% \placeholder{To be completed after the Harmonic Meeting of 30/10: harmonic state-space analysis of the phase-controlled converter over its firing delay, joint Monte Carlo sampling of $L_{ac}$, $L_{dc}$ and $\alpha_{\text{f}}$, the resulting angle distributions and their dependence on the firing delay, and the reference convention of the angles (relative to the terminal phase-$a$ voltage, cosine-referenced).}

% \placeholder{Working assumption, to be replaced: the firing delay has a mean of~$30^{\circ}$ and a Beta-distributed support of~$30^{\circ}$ width, with the shape parameters of Table~\ref{tab:beta_phase_params}, and the overlap angle is independent of the firing delay. The resulting angle model, with one random variable shared by all units of the type, is given in Section~\ref{sec:harmonic_current_angle_uncertainty}.}

% Real networks host a mix of converter technologies, whose harmonic current angles differ in their mean. As a second type, a phase-controlled version of the same 6-pulse converter is considered, i.e., a thyristor rectifier operated at a firing delay~$\alpha_{\text{f}}$. Delaying the firing delays the whole current waveform, which rotates the $\Harmonic$-th harmonic phasor of phase~$a$ by~$-\Harmonic\,\alpha_{\text{f}}$ relative to the uncontrolled converter, while the dispersion caused by the uncertain inductances is retained. The firing delay is set by the operating point, e.g., the demanded dc voltage, and is therefore common to all units of this type and uncertain at the planning stage.

Real networks host a mix of converter technologies. As a second type, the line-commutated converter is equipped with thyristors and operated in phase-controlled mode, resulting in firing delay~$\alpha_{\text{f}}$. Each unit regulates its dc current to a load-dependent reference, assumed common to all units of this type but uncertain at the planning stage from an operator perspective. The reference is expressed as the firing delay~$\alpha_{\text{op}}$ of a unit with the mean commutation inductance~$\bar{L}_{ac}$, so that $\alpha_{\text{f}} = \alpha_{\text{op}}$ for $L_{ac} = \bar{L}_{ac}$. In the absence of operating data, $\alpha_{\text{op}}$ is sampled, as $L_{ac}$ and $L_{dc}$, from a normal distribution with mean of $30^{\circ}$ and standard deviation of $5^{\circ}$, truncated to $[15^{\circ}, 45^{\circ}]$.
% The firing delay at each sample is determined by,
% \begin{equation}
%     \cos\alpha_{\text{f}} = \frac{I_{dc}^{\text{ref}}\left(R_{dc} + \frac{3}{\pi}\omega_1 L_{ac}\right)}{V_{d0}}, 
%     I_{dc}^{\text{ref}} = \frac{V_{d0}\cos\alpha_{\text{op}}}{R_{dc} + \frac{3}{\pi}\omega_1 \bar{L}_{ac}},
% \end{equation}
% where $V_{d0} = 3\sqrt{2}\,V_{LL}/\pi$ and $V_{LL}$ is the line-to-line voltage at ac terminal.
At a given current reference, a larger $L_{ac}$ therefore requires earlier firing to compensate for the increased commutation voltage drop.

Joint Monte Carlo sampling of $L_{ac}$, $L_{dc}$ and $\alpha_{\text{op}}$ ($N = 5000$) yields the harmonic current phase angle distributions for the phase-controlled converter. Following the same uncertainty quantification procedure as for the diode-based converter, the simulated angles are unwrapped about their respective circular means and normalized, and a shared Beta distribution is fitted to the pooled normalized samples. The resulting shared shape parameters, harmonic-specific locations and scales, and fit errors are listed in Table~\ref{tab:beta_phase_params_pc}.
% The histograms from the Monte Carlo simulation and the fitted probability densities are compared in~Fig.~\ref{fig:Shared_Beta_Fit_pc}.
% As summarized in Table~\ref{tab:fund_angle},
The lag of the fundamental current behind the terminal reference voltage~$v_a$ ($v_a = V_a\cos(\omega_1 t - 90^{\circ})$) increases from $9.26^{\circ}$-$16.43^{\circ}$ for the diode-based converter to $20.06^{\circ}$-$46.66^{\circ}$ for the phase-controlled converter. The firing delay~$\alpha_{\text{f}}$ shifts the current waveform later in time and introduces a phase shift of~$-h\alpha_{\text{f}}$ to the $h$-th harmonic. Since uncertainty in operating reference introduces additional order-dependent spread in harmonic phase angles of the phase-controlled converter, the distribution ranges are wider than those of the diode-based converter for every harmonic considered, as indicated in Table~\ref{tab:beta_phase_params} and Table~\ref{tab:beta_phase_params_pc}.

% \begin{table}[!t]
%     \centering
%     \caption{Shared Beta distribution parameters and fit quality for harmonic phase angles of the phase-controlled converter}
%     \label{tab:beta_phase_params_pc}
%     \scriptsize
%     \begin{tabular}{l c c c c}
%     Order & $v_{\min}$ [deg] & $v_{\max}$ [deg] & $v_{\max} - v_{\min}$ [deg] & error \\
%     \hline
%     1st  & $-137.26$ & $-110.15$ & 27.11  & 0.0066 \\
%     5th  & $-141.74$ & $-11.33$  & 130.41 & 0.0013 \\
%     7th  & $-262.58$ & $-49.81$  & 212.77 & 0.0011 \\
%     11th & $-245.62$ & 47.96     & 293.58 & 0.0005 \\
%     13th & $-338.03$ & 20.48     & 358.51 & 0.0007 \\
%     \end{tabular}    
%     \vspace{2pt}
%     \parbox{\columnwidth}{\scriptsize All harmonics share the same shape parameters ($\hat{\alpha}=4.120$ and $\hat{\beta}=3.622$).\\
%     Location parameters $v_{\min}$ and scale parameter $(v_{\max} - v_{\min})$ are harmonic-specific. Angles are unwrapped about their respective circular means.}
% \end{table}



% RMSE compares the 50-bin histogram density with the fitted density evaluated at the bin centers.

% \begin{figure*}[!t]
%     \centering
%     \pgfplotsset{phaseplotpc/.style={shhc, width=0.205\textwidth, height=3cm, ymin=0,
%         xlabel={phase [deg]}, scaled y ticks=false,
%         yticklabel style={/pgf/number format/fixed},
%         every axis plot/.append style={mark=none}},
%         mcspc/.style={ybar interval, fill=gray!30, draw=gray!70, line width=0.2pt},
%         betapc/.style={solid, black, line width=0.9pt}}
%     \subfloat[$h = 1$]{%
%     \begin{tikzpicture}
%         \begin{axis}[phaseplotpc, xtick={-135,-125,-115},
%             ylabel style={align=center}, ylabel={probability\\density\\{[1/deg]}}]
%             \addplot[mcspc] table[x=x, y=density] {journals/shhc/data/type2_shared_beta_phase/hist_interval_h1.dat};
%             \addplot[betapc] table[x=x, y=pdf] {journals/shhc/data/type2_shared_beta_phase/pdf_h1.dat};
%         \end{axis}
%     \end{tikzpicture}}%
%     \subfloat[$h = 5$]{%
%     \begin{tikzpicture}
%         \begin{axis}[phaseplotpc, xtick={-140,-80,-20}]
%             \addplot[mcspc] table[x=x, y=density] {journals/shhc/data/type2_shared_beta_phase/hist_interval_h5.dat};
%             \addplot[betapc] table[x=x, y=pdf] {journals/shhc/data/type2_shared_beta_phase/pdf_h5.dat};
%         \end{axis}
%     \end{tikzpicture}}%
%     \subfloat[$h = 7$]{%
%     \begin{tikzpicture}
%         \begin{axis}[phaseplotpc, xtick={-250,-150,-50}]
%             \addplot[mcspc] table[x=x, y=density] {journals/shhc/data/type2_shared_beta_phase/hist_interval_h7.dat};
%             \addplot[betapc] table[x=x, y=pdf] {journals/shhc/data/type2_shared_beta_phase/pdf_h7.dat};
%         \end{axis}
%     \end{tikzpicture}}%
%     \subfloat[$h = 11$]{%
%     \begin{tikzpicture}
%         \begin{axis}[phaseplotpc, xtick={-200,-100,0}]
%             \addplot[mcspc] table[x=x, y=density] {journals/shhc/data/type2_shared_beta_phase/hist_interval_h11.dat};
%             \addplot[betapc] table[x=x, y=pdf] {journals/shhc/data/type2_shared_beta_phase/pdf_h11.dat};
%         \end{axis}
%     \end{tikzpicture}}%
%     \subfloat[$h = 13$]{%
%     \begin{tikzpicture}
%         \begin{axis}[phaseplotpc, xtick={-300,-150,0}]
%             \addplot[mcspc] table[x=x, y=density] {journals/shhc/data/type2_shared_beta_phase/hist_interval_h13.dat};
%             \addplot[betapc] table[x=x, y=pdf] {journals/shhc/data/type2_shared_beta_phase/pdf_h13.dat};
%         \end{axis}
%     \end{tikzpicture}}
%     \caption{Distribution of harmonic current phase angles of the phase-controlled converter, histograms of 5000 Monte Carlo samples (gray, 50 bins) against the denormalized shared Beta probability density functions (black), with parameters listed in Table~\ref{tab:beta_phase_params_pc}: (a) $h = 1$; (b) $h = 5$; (c) $h = 7$; (d) $h = 11$; (e) $h = 13$.}
%     \label{fig:Shared_Beta_Fit_pc}
% \end{figure*}

% \begin{table}[!t]
%     \centering
%     \caption{The fundamental current angle and displacement of converter types}
%     \label{tab:fund_angle}
%     \scriptsize
%     \begin{tabular}{l c c c}
%     Converter & $\varphi_{1,\min}$ [deg] & $\varphi_{1,\max}$ [deg] & lag behind $v_a$ [deg]  \\
%     \hline
%     Diode-based     & $-106.70$ & $-99.43$  & 9.43--16.70   \\
%     Phase-controlled & $-137.26$ & $-110.15$ & 20.15--47.26  \\
%     \end{tabular}
% \end{table}

% \placeholder{Angle Convention. Can be Removed after Discussion. All phase angles in this section, including those of Tables~\ref{tab:beta_phase_params} and~\ref{tab:beta_phase_params_pc}, are the angles~$\varphi_{\Harmonic}$ of the cosine-based phasors $I_{\Harmonic}\cos(\Harmonic\omega_1 t + \varphi_{\Harmonic})$. The terminal phase-$a$ voltage is set to $v_a = V_a\sin(\omega_1 t) = V_a\cos(\omega_1 t - 90^{\circ})$, so its phasor lies at~$-90^{\circ}$, and the fundamental current lags~$v_a$ by $-(\varphi_1 + 90^{\circ})$. As summarized in Table~\ref{tab:fund_angle},the lag of the fundamental current behind the terminal reference voltage~$v_a$ increases from $9.43^{\circ}$-$16.70^{\circ}$ for the diode-based converter to $20.15^{\circ}$-$47.26^{\circ}$ for the phase-controlled converter.}
\section{Mathematical Framework}
\label{sec:mathematical_framework}

This section presents the mathematical framework for the stochastic harmonic hosting capacity problem of a phase-balanced three-phase power system, considering uncertainty on the angle of the harmonic current injected by the connected harmonic units. To this end, the formulation for the deterministic harmonic hosting capacity presented in~\cite{vanacker_hosting_2026} is extended. For brevity, the original formulation is not repeated in this section, rather the relevant variables and constraints are highlighted and transformed to their stochastic counterpart. %

\subsection{Background on Polynomial Chaos Expansion}
\label{sec:polynomial_chaos_expansion}

This section introduces the necessary background on polynomial chaos expansion, a method for representing random variables in terms of polynomial functions. Concretely, a vector of input random variables, e.g., harmonic current angles, is propagated through the set of constraints of the harmonic hosting capacity problem to obtain a vector of unknown output random variables, e.g., harmonic bus voltages. For a detailed treatment of the subject, the reader is referred to~\cite{sullivan_introduction_2015}.

\subsubsection{Uncertainty Quantification}
%It is assumed that the stochastic input is captured by a real-valued finite-variance stochastic germ~$\bar{\omega} = [\omega_{1},...,\omega_{N}]^{\text{T}}$ of independent random variables. Optimization solvers are unable to handle random variables directly, necessitating a deterministic equivalent of the stochastic problem. For this purpose, polynomial chaos expansion approximates any real-valued random variable~$\RVX$ of finite variance that is a function of the germ~$\bar{\omega}$ as a linear combination~$\ApproxRVX$ of the orthogonal polynomial basis~$[\Polynoom_{\PCEIndex}(\bar{\omega})]_{\PCEIndex \in \PCESet}$, where $\PCESet\subseteq\NaturalNumberSet_{0}$, as,
It is assumed that the stochastic input is captured by a real-valued finite-variance stochastic germ~$\bar{\omega} = [\omega_{1},...,\omega_{N}]^{\text{T}}$ of independent random variables\changed{, i.e., the inputs from which all random quantities of the problem are derived}. Optimization solvers are unable to handle random variables directly, necessitating a deterministic equivalent of the stochastic problem. For this purpose, polynomial chaos expansion approximates any real-valued random variable~$\RVX$ of finite variance that is a function of the germ~$\bar{\omega}$ as a linear combination~$\ApproxRVX$ of the orthogonal polynomial basis~$[\Polynoom_{\PCEIndex}(\bar{\omega})]_{\PCEIndex \in \PCESet}$, where $\PCESet\subseteq\NaturalNumberSet_{0}$, as,
\begin{IEEEeqnarray}{ r C l }
    \ApproxRVX &=& \sum_{\PCEIndex \in \PCESet} \CoefficientX_{\PCEIndex} \cdot \Polynoom_{\PCEIndex} =
    \sum_{\PCEIndex \in \PCESet} \frac{\langle \RVX, \Polynoom_{\PCEIndex} \rangle}{\langle \Polynoom_{\PCEIndex}, \Polynoom_{\PCEIndex} \rangle} \cdot \Polynoom_{\PCEIndex}, \label{eq:gpc_expansion}
\end{IEEEeqnarray}
where~$\CoefficientX_{\PCEIndex} \in \RealNumberSet$ denote the polynomial chaos expansion coefficients and $\langle \cdot,\cdot \rangle$ denotes the inner product with respect to the probability measure of the germ. Following the Askey scheme~\cite{sullivan_introduction_2015}, each Beta-distributed random variable corresponds to the Jacobi family of orthogonal polynomials; the multivariate basis is the tensor product of the univariate ones, truncated to total degree~$\delta$, and is constructed numerically with \textsc{PolyChaos.jl}~\cite{muhlpfordt_polychaos_2020}.

\subsubsection{Uncertainty Projection onto Equality Constraints}
Through polynomial chaos expansion, any stochastic equality constraint is equivalently represented by $|\PCESet|$ deterministic equality constraints. The cardinality of the model~$|\PCESet|$ is determined by the dimension of the germ and the highest polynomial degree as,
\begin{IEEEeqnarray}{ r C l }
    |\PCESet| &=& \frac{(|\bar{\omega}| + \delta)!}{|\bar{\omega}|!\,\delta!}, \label{eq:cardinality}
\end{IEEEeqnarray}
where~$\delta$ denotes the maximum degree of the polynomial basis functions.

The Galerkin projection forms a system of equations for the stochastic modes~(Chapter~12 in~\cite{sullivan_introduction_2015}). For the harmonic hosting capacity problem, two operations stand out: summation and multiplication. For three approximated random variables~$\ApproxRVX$,~$\ApproxRVY$ and~$\ApproxRVZ$ defined on the same basis, the Galerkin projection, denoted by~$\gp$, of both operations is,
\begin{IEEEeqnarray}{ l C l l}
    \ApproxRVZ = \ApproxRVX + \ApproxRVY &\gp& \CoefficientZ_{\PCEIndex} = \CoefficientX_{\PCEIndex} + \CoefficientY_{\PCEIndex},~&\forall \PCEIndex \in \PCESet,~ \label{eq:gp_sum} \\
    \ApproxRVZ = \ApproxRVX \cdot \ApproxRVY &\gp& \CoefficientZ_{\PCEIndex} = \mkern-15mu \sum_{\bar{\kappa} \in P^{*}(\PCESet,2)} \mkern-15mu M_{\bar{\kappa},\PCEIndex} \,(\CoefficientX_{\bar{\kappa}_{1}} \cdot \CoefficientY_{\bar{\kappa}_{2}}),\quad&\forall \PCEIndex \in \PCESet,~ \label{eq:gp_product}
\end{IEEEeqnarray}
where~$P^{*}(\PCESet,2)$ denotes the 2-tuples of~$\PCESet$ and the multiplication tensor,
\begin{IEEEeqnarray}{ r C l }
    M_{\bar{\kappa},\PCEIndex} = \frac{\langle \Polynoom_{\bar{\kappa}_{1}}\Polynoom_{\bar{\kappa}_{2}},\Polynoom_{k} \rangle}{\langle \Polynoom_{k},\Polynoom_{k}\rangle}, \label{eq_m_tensor_def}
\end{IEEEeqnarray}
is computed ahead of time. For $\PCEIndex=0$, orthogonality implies that $M_{\bar{\kappa},0}$ is diagonal with strictly positive entries; for $\PCEIndex \neq 0$, the tensor has mixed-sign entries, and the matrix $[M_{\bar{\kappa},k}]_{\bar\kappa}$ is generally indefinite.

\subsubsection{Chance-Constraint Reformulation of Inequality Constraints}
All harmonic limits in this work are upper bounds. A chance constraint ensures that the probability of~$\RVX$ exceeding a deterministic limit~$x^{\text{max}}$ is below a risk level~$\varepsilon$, i.e., $\Probability(\RVX \leq x^{\text{max}}) \geq 1 - \varepsilon$. It is enforced through the moment-based reformulation,
\begin{IEEEeqnarray}{ r C l }
    \Expectation(\RVX) + \lambda(\varepsilon)\sqrt{\Variance(\RVX)} \leq x^{\text{max}}, \label{eq:cc_constraint}
\end{IEEEeqnarray}
with the expectation and variance of an approximated random variable~$\ApproxRVX$ given in terms of its coefficients as,
\begin{IEEEeqnarray}{ r C l }
    \Expectation[\ApproxRVX]   &=& \CoefficientX_{0}, \label{eq:gpc_expectation} \\
    \Variance[\ApproxRVX]      &=& \mkern-10mu \sum_{\PCEIndex \in \PCESet \backslash \{0\}} \mkern-10mu \langle \Polynoom_{\PCEIndex}, \Polynoom_{\PCEIndex} \rangle \, \CoefficientX_{\PCEIndex}^{2}, \label{eq:gpc_variance}
\end{IEEEeqnarray}
so that~\eqref{eq:cc_constraint} is a second-order cone constraint in the coefficients~$\{\CoefficientX_{\PCEIndex}\}_{\PCEIndex \in \PCESet}$. The risk factor~$\lambda(\varepsilon)$ is calibrated per constraint as $\lambda^{*} = (q_{1-\varepsilon}(\RVX)-\Expectation[\ApproxRVX])/\sqrt{\Variance[\ApproxRVX]}$, with the $(1-\varepsilon)$ quantile~$q_{1-\varepsilon}$ of samples of the constrained quantity and the moments~\eqref{eq:gpc_expectation}--\eqref{eq:gpc_variance} of its expansion, and is updated with the hosting capacity until convergence. As it is calibrated on the moments of the expansion, the bound equals the sampled quantile irrespective of the truncation of the basis, i.e., it targets the risk level~$\varepsilon$ exactly.

\subsection{Uncertainty on Angle of Harmonic Current Injection}
\label{sec:harmonic_current_angle_uncertainty}

%This section discusses the treatment of the uncertain angles~$\phi_{\Unit,\Harmonic}$ of the harmonic currents injected by the harmonic units~$\Unit \in \HarmonicUnitSet$, for all harmonics~$\Harmonic \in \HarmonicSet$. Throughout, a polynomial basis of degree~$\delta = 2$ is used, which captures the non-linear nature of the power flow equations without significant loss of accuracy~\cite{muhlpfordt_chance_2019}. For narrow angle ranges, it keeps at least~92\,\% of the variance of the cosine and sine of the angle up to the 13th harmonic, but only~39\,\% at the 25th harmonic of the thyristor bridges. The calibration of the risk factor on the moments of the expansion, Section~\ref{sec:polynomial_chaos_expansion}, compensates for this truncation.
This section discusses the treatment of the uncertain angles~$\phi_{\Unit,\Harmonic}$ of the harmonic currents injected by the harmonic units~$\Unit \in \HarmonicUnitSet$, for all harmonics~$\Harmonic \in \HarmonicSet$. Throughout, a polynomial basis of degree~$\delta = 2$ is used, which captures the non-linear nature of the power flow equations without significant loss of accuracy~\cite{muhlpfordt_chance_2019}. \changed{For the angle distributions of Section~\ref{sec:uncertainty}, computed by quadrature,} it keeps at least~92\,\% of the variance of the cosine and sine of the angle up to the 13th harmonic, but only~39\,\% at the 25th harmonic of the thyristor bridges. The calibration of the risk factor on the moments of the expansion, Section~\ref{sec:polynomial_chaos_expansion}, compensates for this truncation.

\subsubsection{Random Variable per Harmonic Unit}
As shown in Section~\ref{sec:uncertainty}, the phase angles of a converter across the harmonic orders are strongly dependent, and their marginal distributions are well approximated by a single shared Beta distribution with harmonic-specific location and scale. Each harmonic unit~$\Unit$ is therefore assigned one random variable~$\omega_{\Unit} \sim \mathcal{B}(\hat{\alpha},\hat{\beta})$ on~$[0,1]$, which drives its angle at every harmonic through an affine map,
\begin{IEEEeqnarray}{ r C l }
    \phi_{\Unit,\Harmonic} &=& v^{\text{min}}_{\Harmonic} + (v^{\text{max}}_{\Harmonic} - v^{\text{min}}_{\Harmonic})\,\omega_{\Unit} + \Harmonic\,\theta_{\Node(\Unit)}, \label{eq:angle_affine}
\end{IEEEeqnarray}
while the random variables of different units are independent. At the higher harmonics, the angle range may exceed~$2\pi$, e.g., 660$^{\circ}$ at the 25th harmonic of the thyristor bridges. The affine map then wraps around the circle, which leaves the phasor~$e^{\im\phi_{\Unit,\Harmonic}}$ and the dependence between the harmonics of a unit well defined. For a mix of converter technologies, the shape parameters~$\hat{\alpha}$, $\hat{\beta}$ and the angle ranges~$[v^{\text{min}}_{\Harmonic}, v^{\text{max}}_{\Harmonic}]$ are those of the technology of unit~$\Unit$, Section~\ref{sec:uncertainty}, and the basis is the tensor product of the univariate bases of the units.

\subsubsection{Local Angle Reference}
The angle ranges of Section~\ref{sec:uncertainty} result from a converter model driven by its own terminal voltage, i.e., they are defined relative to the fundamental voltage at the bus~$\Node(\Unit)$ the unit connects to. A shift of the fundamental by its angle~$\theta_{\Node}$ rotates the $\Harmonic$-th harmonic phasor by~$\Harmonic\,\theta_{\Node}$, which gives the last term of~\eqref{eq:angle_affine}. The fundamental angles follow from the fundamental power flow, which is solved first. Since~$\theta_{\Node}$ depends on the operating point and not only on the network topology, units of the same type at different buses inject at different absolute angles: the local angle reference is a physically required source of angle diversity, and it is the reference used throughout this work.

\subsubsection{Exact Projection of the Trigonometric Functions}
In the deterministic formulation~\cite{vanacker_hosting_2026}, the harmonic current magnitude~$\CurrentMagnitude_{\Unit,\Harmonic}$ is related to its rectangular components through the angle~$\phi_{\Unit,\Harmonic}$,
\begin{IEEEeqnarray}{ l }
        \cos(\phi_{\Unit,\Harmonic}) \, \CurrentMagnitude_{\Unit,\Harmonic} =
        \RealCurrent_{\Unit,\Harmonic}, \quad \forall \Unit \in \HarmonicUnitSet, \Harmonic \in \HarmonicSet,\label{eq:real_harmonic_unit_current} \\
        \sin(\phi_{\Unit,\Harmonic}) \, \CurrentMagnitude_{\Unit,\Harmonic} =
        \ImaginaryCurrent_{\Unit,\Harmonic}, \quad \forall \Unit \in \HarmonicUnitSet, \Harmonic \in \HarmonicSet. \label{eq:imag_harmonic_unit_current}
\end{IEEEeqnarray}
Since the angle~\eqref{eq:angle_affine} is affine in a single random variable, $\cos(\phi_{\Unit,\Harmonic})$ and $\sin(\phi_{\Unit,\Harmonic})$ are functions of~$\omega_{\Unit}$ alone, and their coefficients follow from~\eqref{eq:gpc_expansion}, for every index~$\PCEIndex$ supported on~$\omega_{\Unit}$ alone and zero otherwise,
\begin{IEEEeqnarray}{ r C l }
    \cos(\phi_{\Unit,\Harmonic})_{\PCEIndex} &=& \frac{\langle \cos(\phi_{\Unit,\Harmonic}), \Polynoom_{\PCEIndex} \rangle}{\langle \Polynoom_{\PCEIndex},\Polynoom_{\PCEIndex}\rangle}, \label{eq:trig_quadrature}
\end{IEEEeqnarray}
and analogously for the sine. %
%The coefficients~$\cos(\phi_{\Unit,\Harmonic})_{\PCEIndex}$ and~$\sin(\phi_{\Unit,\Harmonic})_{\PCEIndex}$ are parameters, not optimization variables.
The coefficients~$\cos(\phi_{\Unit,\Harmonic})_{\PCEIndex}$ and~$\sin(\phi_{\Unit,\Harmonic})_{\PCEIndex}$ are parameters, not optimization variables\changed{, and therefore do not affect the convexity of the optimization problem}.

\subsection{General Stochastic Harmonic Hosting Capacity}
\label{sec:stochastic_harmonic_hosting_capacity}

The stochastic model determines a deterministic harmonic hosting capacity given uncertainty on the harmonic current angles, while complying with chance constraints on (harmonic) voltages and currents. Given that the problem is formulated from the perspective of the power system operator, deterministic harmonic budgets are the practical outcome for its customers. A deterministic hosting capacity implies that~$\CurrentMagnitude_{\Unit,\Harmonic,\PCEIndex} = 0$ for any non-zero index~$\PCEIndex$, so the counterparts of~\eqref{eq:real_harmonic_unit_current} and~\eqref{eq:imag_harmonic_unit_current} are,
\begin{IEEEeqnarray}{ l }
    \cos(\phi_{\Unit,\Harmonic})_{\PCEIndex} \, \CurrentMagnitude_{\Unit,\Harmonic,0} = \RealCurrent_{\Unit,\Harmonic,\PCEIndex},
        \quad \forall \Unit \in \HarmonicUnitSet, \Harmonic \in \HarmonicSet, \PCEIndex \in \PCESet,
    \label{eq:real_harmonic_unit_current_pce} \\
    \sin(\phi_{\Unit,\Harmonic})_{\PCEIndex} \, \CurrentMagnitude_{\Unit,\Harmonic,0} = \ImaginaryCurrent_{\Unit,\Harmonic,\PCEIndex},
        \quad \forall \Unit \in \HarmonicUnitSet, \Harmonic \in \HarmonicSet, \PCEIndex \in \PCESet,
    \label{eq:imag_harmonic_unit_current_pce}
\end{IEEEeqnarray}
which are linear in the decision variables. Any variables related to the fairness objective are also deterministic. All other parameters, e.g., network impedances, are deterministic.

\subsubsection{Reformulation of Linear Constraints}
\label{sec:reformulation_of_linear_constraints}
The network equations of the current-voltage formulation in~\cite{vanacker_hosting_2026}, i.e., the reference bus, Kirchhoff's current law, Ohm's law, branch current and transformer constraints~(10)--(31), are linear, as are the fairness constraints~(34), (36) and~(38)--(39) and the harmonic current injection constraints~(40)--(41), of which the latter are treated above. Supported by~\eqref{eq:gp_sum}, a generic linear constraint $a\,\RVX + b\,\RVY = c$ with deterministic parameters is repeated over all indices, $a\,\RVX_{k} + b\,\RVY_{k} = c_{k}$, $\forall \PCEIndex \in \PCESet$, where $c_{k} = c$ for $k=0$ and $c_{k}=0$ otherwise. Since each voltage and current coefficient is then linear in the injected current coefficients, the physical variables are non-zero only on the indices supported by a single random variable%
, which reduces the number of variables.

\subsubsection{Reformulation of Quadratic Constraints and Problem Class}
\label{sec:reformulation_of_quadratic_constraints}
The root-mean-square, total and individual harmonic distortion voltage limits and the root-mean-square current limit are the second-order cone constraints~(47)--(49) and~(51) in~\cite{vanacker_hosting_2026}. Consider a generic constraint $\RVX^{2} + \RVY^{2} \leq a^{2}$ with deterministic~$a$. Following~\cite{vanacker_ivr}, an auxiliary lifted variable~$\RVZ$ is introduced, $\RVX^{2} + \RVY^{2} = \RVZ \leq a^{2}$. Supported by~\eqref{eq:gp_product}, the lifting yields~$|\PCESet|$ constraints,
\begin{IEEEeqnarray}{ l }
    \sum_{\bar{\kappa} \in P^{*}(\PCESet,2)} \mkern-15mu M_{\bar{\kappa},\PCEIndex} \,(\RVX_{\kappa_{1}} \RVX_{\kappa_{2}} + \RVY_{\kappa_{1}} \RVY_{\kappa_{2}}) = \RVZ_{\PCEIndex}, \quad \forall \PCEIndex \in \PCESet. \label{eq:soc_aux_eq}
\end{IEEEeqnarray}
For~$\PCEIndex = 0$ the quadratic form is positive semidefinite, but for~$\PCEIndex \neq 0$ it is indefinite and the equality must be retained: relaxing it would allow~$\RVZ_k = 0$ where the true coefficient is negative, which underestimates the variance and makes the chance constraint over-optimistic. The right-hand side becomes the chance constraint,
\begin{IEEEeqnarray}{ l }
    \RVZ_{0} + \lambda(\varepsilon) \sqrt{ \sum_{\PCEIndex \in \PCESet \backslash \{0\}} \mkern-10mu \langle \Polynoom_{\PCEIndex}, \Polynoom_{\PCEIndex} \rangle \, \RVZ_{\PCEIndex}^{2}} \leq a^{2}, \label{eq:cc_soc}
\end{IEEEeqnarray}
a second-order cone in~$\{\RVZ_{\PCEIndex}\}_{\PCEIndex\in\PCESet}$. With the equalities~\eqref{eq:soc_aux_eq} for~$\PCEIndex \neq 0$, the general stochastic harmonic hosting capacity problem is a non-convex quadratically constrained quadratic program.

\subsection{Reduction to a Second-Order Cone Program}
\label{sec:absolute_equality_socp}

The general problem reduces to a second-order cone program for every fairness criterion that fixes the ratios between the allocations of the harmonic units, i.e., $\CurrentMagnitude_{\Unit,\Harmonic,0} = w_{\Unit}\,\HostingCapacity_{\Harmonic}$ with fixed positive weights~$w_{\Unit}$ and a single hosting capacity~$\HostingCapacity_{\Harmonic}$ per harmonic. Absolute equality, $w_{\Unit} = 1$, and weighted equality, e.g., with weights proportional to the rated power of the units, are such criteria. The Kalai-Smorodinsky solution~(37)--(39) in~\cite{vanacker_hosting_2026} is one when all ratios to the individual maxima are equal at the optimum, with~$w_{\Unit}$ the individual maximum of unit~$\Unit$, and maximin~(35) in~\cite{vanacker_hosting_2026} coincides with absolute equality when raising the current of a unit cannot relax a binding constraint. Criteria with a free allocation, e.g., maximum efficiency or proportional fairness~\cite{chen_guide_2023}, require a current per unit and lead to a semidefinite relaxation instead.

For fixed weights, \eqref{eq:real_harmonic_unit_current_pce}--\eqref{eq:imag_harmonic_unit_current_pce} express every injected current coefficient as a known parameter times~$\HostingCapacity_{\Harmonic}$. As the network equations are linear in the coefficients, all voltage coefficients are linear in~$\HostingCapacity_{\Harmonic}$ as well,
\begin{IEEEeqnarray}{ l }
    \RealVoltage_{\Node,\Harmonic,\PCEIndex} = C^{\text{re}}_{\Node,\Harmonic,\PCEIndex}\, \HostingCapacity_{\Harmonic}, \qquad
    \ImaginaryVoltage_{\Node,\Harmonic,\PCEIndex} = C^{\text{im}}_{\Node,\Harmonic,\PCEIndex}\, \HostingCapacity_{\Harmonic},
    \label{eq:voltage_linear}
\end{IEEEeqnarray}
where~$C^{\text{re}}$ and~$C^{\text{im}}$ follow from the harmonic transfer impedances~$Z_{\Harmonic}$ of the network, obtained by inverting the harmonic admittance matrix, the weights~$w_{\Unit}$ and the trigonometric coefficients~\eqref{eq:trig_quadrature}; branch currents follow analogously. The transformer vector-group enters the admittance matrix with the sign of the sequence of the harmonic. Zero-sequence harmonics additionally require the winding configuration and earthing of the transformers. The lifting then collapses to the precomputed scalars,
\begin{IEEEeqnarray}{ l }
\alpha_{\Node,\Harmonic,\PCEIndex} = \mkern-18mu \sum_{\bar{\kappa} \in P^{*}(\PCESet,2)} \mkern-18mu M_{\bar{\kappa},\PCEIndex} \big( C^{\text{re}}_{\Node,\Harmonic,\bar{\kappa}_{1}} C^{\text{re}}_{\Node,\Harmonic,\bar{\kappa}_{2}} + C^{\text{im}}_{\Node,\Harmonic,\bar{\kappa}_{1}} C^{\text{im}}_{\Node,\Harmonic,\bar{\kappa}_{2}} \big),
\end{IEEEeqnarray}
i.e., the coefficients of~$|\ComplexVoltage_{\Node,\Harmonic}|^2/\HostingCapacity_{\Harmonic}^2$. With~$\LiftedHostingCapacity_{\Harmonic} \geq \HostingCapacity_{\Harmonic}^{2}$, a rotated second-order cone, every chance constraint on a sum of squared harmonic quantities over a set~$\mathcal{H}_{c} \subseteq \HarmonicSet$ of harmonics becomes,
\begin{IEEEeqnarray}{ l }
    \sum_{\Harmonic \in \mathcal{H}_{c}} \alpha_{\Node,\Harmonic,0}\, \LiftedHostingCapacity_{\Harmonic} \leq a_{\Node,c} \nonumber \\
    \quad -\: \lambda(\varepsilon) \sqrt{ \sum_{\PCEIndex \in \PCESet \backslash \{0\}} \mkern-10mu \langle \Polynoom_{\PCEIndex}, \Polynoom_{\PCEIndex} \rangle \Big( \sum_{\Harmonic \in \mathcal{H}_{c}} \alpha_{\Node,\Harmonic,\PCEIndex}\, \LiftedHostingCapacity_{\Harmonic} \Big)^{2}}, \label{eq:socp_cc}
\end{IEEEeqnarray}
a second-order cone in~$\bar{\LiftedHostingCapacity}$: $\mathcal{H}_{c} = \{\Harmonic\}$ for the individual harmonic voltage distortion, which reduces to a linear bound on~$\LiftedHostingCapacity_{\Harmonic}$, and $\mathcal{H}_{c} = \HarmonicSet$ for the total harmonic distortion and the root-mean-square voltage and current limits, where~$a_{\Node,c}$ is the limit net of the fundamental. The stochastic harmonic hosting capacity is thus the second-order cone program,
\begin{IEEEeqnarray}{ l l l }
    \max_{\bar{\HostingCapacity},\,\bar{\LiftedHostingCapacity}} & \sum_{\Harmonic \in \HarmonicSet} \Big( \sum_{\Unit \in \HarmonicUnitSet} w_{\Unit} \Big) \HostingCapacity_{\Harmonic} & \nonumber \\
    \text{subject to} \quad & \HostingCapacity_{\Harmonic}^{2} \leq \LiftedHostingCapacity_{\Harmonic}, & \forall \Harmonic \in \HarmonicSet, \nonumber \\
    & w_{\Unit}\,\HostingCapacity_{\Harmonic} \leq \RatedCurrent_{\Unit}, & \forall \Unit \in \HarmonicUnitSet, \Harmonic \in \HarmonicSet, \nonumber \\
    & w_{\Unit}^{2}\,\LiftedHostingCapacity_{\Harmonic} \leq (\RatedCurrent_{\Unit})^{2}, \quad & \forall \Unit \in \HarmonicUnitSet, \Harmonic \in \HarmonicSet, \nonumber \\
    & \eqref{eq:socp_cc}, & \label{eq:socp}
\end{IEEEeqnarray}
with one pair of variables per harmonic, whatever the size of the network or the cardinality of the basis; the latter only enters the precomputed scalars~$\alpha$. The relaxation~$\HostingCapacity_{\Harmonic}^{2} \leq \LiftedHostingCapacity_{\Harmonic}$ is exact: the chance constraints~\eqref{eq:socp_cc} involve~$\bar{\LiftedHostingCapacity}$ only and the objective is strictly increasing in each~$\HostingCapacity_{\Harmonic}$, so any optimum satisfies~$\HostingCapacity_{\Harmonic} = \sqrt{\LiftedHostingCapacity_{\Harmonic}}$. The rating is therefore imposed on~$\LiftedHostingCapacity_{\Harmonic}$ as well. Since the left-hand side of~\eqref{eq:socp_cc} is not monotone in~$\bar{\LiftedHostingCapacity}$, i.e., under recombination an increase of~$\LiftedHostingCapacity_{\Harmonic}$ may reduce the variance more than it increases the expectation, a bound on~$\HostingCapacity_{\Harmonic}$ alone would allow~$\LiftedHostingCapacity_{\Harmonic} > \HostingCapacity_{\Harmonic}^{2}$ once the rating binds, and the reported hosting capacity would violate the chance constraints. A limit that the fundamental already exhausts, $a_{\Node,c} \leq 0$, leaves no harmonic headroom and forces the contributing~$\LiftedHostingCapacity_{\Harmonic}$ to zero. When every angle range collapses to a point, all variances vanish and~\eqref{eq:socp} reduces exactly to the deterministic model.


\input{journals/shhc/sections/numerical_illustration}
\input{journals/shhc/sections/conclusion}

% \input{journals/shhc/sections/supplementary}

\section*{Acknowledgement}
The authors used Claude (Anthropic) to assist with implementing and testing parts of the open-source code in \textsc{HarmonicPowerModels.jl}, running the numerical studies of Section~\ref{sec:numerical_illustration}, generating the TikZ figures from the resulting data, and drafting and editing the text of Sections~\ref{sec:introduction}, \ref{sec:mathematical_framework}, \ref{sec:numerical_illustration} and~\ref{sec:conclusion}. All AI-assisted content was reviewed, verified and revised by the authors, who take full responsibility for the content of this article.

The authors thank dr. Kaan Yurtseven for the fruitful discussions and dr. Johan~Rimez for proof reading the paper.

\bibliographystyle{IEEEtran}
\bibliography{journals/shhc/shhc_references}


\end{document}